% =============================================================================
%  fall_n — Technical Documentation
% =============================================================================
%
%  Compile:
%    pdflatex main && bibtex main && pdflatex main && pdflatex main
%
% =============================================================================
\documentclass[11pt,a4paper]{report}

% ─── Packages ────────────────────────────────────────────────────────────
\usepackage[utf8]{inputenc}
\usepackage[T1]{fontenc}
\usepackage[english]{babel}
\usepackage{lmodern}

\usepackage{amsmath,amssymb,amsthm,mathtools}
\usepackage{bm}
\usepackage{booktabs}
\usepackage{array}
\usepackage{graphicx}
\usepackage{xcolor}
\usepackage{listings}
\usepackage{hyperref}
\usepackage{cleveref}
\usepackage{algorithm2e}
\usepackage{geometry}
\usepackage{enumitem}
\usepackage{caption}
\usepackage{subcaption}
\usepackage{fancyhdr}
\usepackage{tikz}
\usetikzlibrary{calc,positioning,arrows.meta}

\geometry{margin=2.5cm}

% ─── Listings style ──────────────────────────────────────────────────────
\lstdefinestyle{cppstyle}{
  language=C++,
  basicstyle=\ttfamily\small,
  keywordstyle=\bfseries\color{blue!70!black},
  commentstyle=\itshape\color{green!50!black},
  stringstyle=\color{red!60!black},
  numbers=left,
  numberstyle=\tiny\color{gray},
  breaklines=true,
  frame=single,
  framerule=0.4pt,
  rulecolor=\color{gray!50},
  tabsize=4,
  showstringspaces=false,
  morekeywords={constexpr, consteval, noexcept, decltype, auto, nullptr,
                requires, concept, std, size_t, span}
}
\lstset{style=cppstyle}

% ─── Math operators ──────────────────────────────────────────────────────
\DeclareMathOperator{\tr}{tr}
\DeclareMathOperator{\dev}{dev}
\DeclareMathOperator{\sym}{sym}
\DeclareMathOperator{\diag}{diag}
\newcommand{\R}{\mathbb{R}}
\newcommand{\N}{\mathbb{N}}
\newcommand{\dd}{\mathrm{d}}
\newcommand{\bxi}{\bm{\xi}}
\newcommand{\blambda}{\bm{\lambda}}
\newcommand{\bsigma}{\bm{\sigma}}
\newcommand{\bvarepsilon}{\bm{\varepsilon}}
\newcommand{\bF}{\mathbf{F}}
\newcommand{\bC}{\mathbf{C}}
\newcommand{\bS}{\mathbf{S}}
\newcommand{\bP}{\mathbf{P}}
\newcommand{\bI}{\mathbf{I}}
\newcommand{\calS}{\mathcal{S}}
\newcommand{\calT}{\mathcal{T}}

% ─── Theorem environments ────────────────────────────────────────────────
\theoremstyle{definition}
\newtheorem{definition}{Definition}[chapter]
\newtheorem{remark}{Remark}[chapter]
\newtheorem{proposition}{Proposition}[chapter]

% ─── Document info ───────────────────────────────────────────────────────
\title{%
  \textbf{\texttt{fall\_n}} --- Technical Documentation\\[1ex]
  \large Simplex Elements, Quadrature Theory,\\
  Element Architecture, and Implementation Decisions}
\author{S. E. Chavarría-Miranda}
\date{\today}

% ═════════════════════════════════════════════════════════════════════════
\begin{document}

\maketitle

\begin{abstract}
This document provides a comprehensive technical reference for the
\texttt{fall\_n} nonlinear finite element library.  It covers the
implementation of simplex (tetrahedral) finite elements, numerical integration
strategies on simplicial domains, the engineering decisions made to ensure
robust post-processing, and the layered element architecture with its
policy-based design patterns.

Key topics include:
\begin{itemize}[nosep]
  \item Lagrange shape functions on the reference simplex and their isoparametric
        mapping to physical space;
  \item Classical simplex quadrature rules and their negative-weight pathology;
  \item The Stroud conical product rule with Gauss--Jacobi nodes via the
        Golub--Welsch algorithm as an all-positive-weight alternative;
  \item Diagnosis and resolution of the checkerboard artefact in lumped $L^2$
        projections on quadratic tetrahedra;
  \item Volume-weighted (SPR-type) nodal averaging as a robust projection
        method;
  \item A bug fix in PETSc~3.24.5 \texttt{section.c} and application-level
        workarounds for PETSc indexing and preallocation issues;
  \item The element architecture design: External Polymorphism at the geometry
        and assembly levels, policy-based customisation points
        (MaterialPolicy, KinematicPolicy, ElementPolicy), compile-time
        efficiency, and comparison with OpenSees and Kratos Multiphysics.
\end{itemize}
Each chapter includes the relevant mathematical theory with bibliographic
references, the concrete C++ implementation in \texttt{fall\_n}, and a
discussion of the design rationale.
\end{abstract}

\tableofcontents

% ─── Chapters ────────────────────────────────────────────────────────────
% =============================================================================
%  Chapter 1 — Simplex Finite Elements
% =============================================================================

\chapter{Simplex Finite Elements}
\label{ch:simplex}

This chapter describes the construction of Lagrangian finite elements on
simplicial cells (triangles in 2D, tetrahedra in 3D) as implemented in
\texttt{fall\_n}.  The development follows the classical Ciarlet
definition~\cite{Ciarlet1972} specialised to Lagrange interpolation on the
reference simplex.


% ─────────────────────────────────────────────────────────────────────────
\section{Reference Simplex}
\label{sec:refsimplex}

\begin{definition}[Reference simplex]
The \emph{reference simplex} in $\R^D$ is the convex hull
\begin{equation}
  \calS_D \;=\; \Bigl\{
    \bxi = (\xi_1, \dots, \xi_D) \in \R^D
    \;\Big|\;
    \xi_i \ge 0,\;
    \sum_{i=1}^{D} \xi_i \le 1
  \Bigr\},
  \label{eq:ref_simplex}
\end{equation}
with \emph{vertices}
\[
  \mathbf{v}_0 = \mathbf{0}, \qquad
  \mathbf{v}_i = \mathbf{e}_i \quad (i = 1, \dots, D),
\]
where $\mathbf{e}_i$ is the $i$-th standard basis vector.
\end{definition}

The \emph{barycentric coordinates} $(\lambda_0, \lambda_1, \dots, \lambda_D)$
are defined as
\begin{equation}
  \lambda_0 = 1 - \sum_{i=1}^{D} \xi_i, \qquad
  \lambda_i = \xi_i \quad (i = 1, \dots, D).
  \label{eq:bary}
\end{equation}
They form a partition of unity, $\sum_{i=0}^{D} \lambda_i = 1$, and satisfy
$\lambda_i \ge 0$ on $\calS_D$.

The volume of the reference simplex is
\begin{equation}
  |\calS_D| = \frac{1}{D!}\,.
  \label{eq:simplex_volume}
\end{equation}


% ─────────────────────────────────────────────────────────────────────────
\section{Lagrange Shape Functions}
\label{sec:shape_functions}

\subsection{Linear Elements (Order 1)}

The linear simplex has $D+1$ nodes (the vertices).  The shape functions are
simply the barycentric coordinates:
\begin{equation}
  N_i(\bxi) = \lambda_i(\bxi), \qquad i = 0, 1, \dots, D.
  \label{eq:shape_linear}
\end{equation}

Since $\lambda_i \ge 0$ on $\calS_D$ and $\lambda_i(\mathbf{v}_j) =
\delta_{ij}$, these are the unique linear polynomials that interpolate the
nodal values.

\paragraph{Concrete instances in \texttt{fall\_n}:}
\begin{itemize}[nosep]
  \item TRI3 ($D=2$, 3 nodes): \texttt{SimplexElement<3,2,1>} (triangle in 3D
        ambient space, used for boundary faces)
  \item TET4 ($D=3$, 4 nodes): \texttt{SimplexElement<3,3,1>}
\end{itemize}


\subsection{Quadratic Elements (Order 2)}
\label{sec:quadratic_simplex}

The quadratic simplex has
\begin{equation}
  n_{\text{nodes}} = \binom{D+2}{2} = \frac{(D+1)(D+2)}{2}
  \label{eq:numnodes_quadratic}
\end{equation}
nodes: the $D+1$ vertices plus $\binom{D+1}{2}$ edge midpoints.

\paragraph{Vertex shape functions.}
For vertex $i$ ($i=0,\dots,D$):
\begin{equation}
  \boxed{N_i(\bxi) = \lambda_i\,(2\lambda_i - 1)}
  \label{eq:shape_vertex_quad}
\end{equation}

\begin{remark}[Sign of vertex shape functions]
\label{rem:vertex_negative}
The vertex shape function $N_i = \lambda_i(2\lambda_i - 1)$ is \emph{negative}
whenever $0 < \lambda_i < \tfrac{1}{2}$.  This region is non-empty in the
interior of $\calS_D$.  In particular, interior quadrature points that are
``far'' from vertex $i$ will have $\lambda_i \ll \tfrac{1}{2}$, making
$N_i < 0$ at those points.

This fact is benign for stiffness matrix assembly (the bilinear form involves
gradients, not function values).  However, it becomes critical for lumped $L^2$
projection, as discussed in~\cref{ch:negative_weights}.
\end{remark}

\paragraph{Edge midpoint shape functions.}
For the midpoint of edge $(i, j)$ with $0 \le i < j \le D$:
\begin{equation}
  \boxed{N_{ij}(\bxi) = 4\,\lambda_i\,\lambda_j}
  \label{eq:shape_edge_quad}
\end{equation}

These are always non-negative on $\calS_D$, since both $\lambda_i, \lambda_j
\ge 0$.

\paragraph{Concrete instances:}
\begin{itemize}[nosep]
  \item TRI6 ($D=2$, 6 nodes): \texttt{SimplexElement<3,2,2>}
  \item TET10 ($D=3$, 10 nodes): \texttt{SimplexElement<3,3,2>}
\end{itemize}


% ─────────────────────────────────────────────────────────────────────────
\section{Shape Function Derivatives}
\label{sec:derivatives}

Since $\lambda_0 = 1 - \xi_1 - \cdots - \xi_D$ and $\lambda_i = \xi_i$ for
$i \ge 1$, the barycentric coordinate gradients (with respect to $\bxi$) are:
\begin{equation}
  \frac{\partial \lambda_0}{\partial \xi_k} = -1, \qquad
  \frac{\partial \lambda_i}{\partial \xi_k} = \delta_{ik}
  \quad (i = 1, \dots, D).
  \label{eq:bary_grad}
\end{equation}

For the quadratic vertex shape function $N_i = \lambda_i(2\lambda_i - 1)$:
\begin{equation}
  \frac{\partial N_i}{\partial \xi_k}
  = (4\lambda_i - 1)\,\frac{\partial \lambda_i}{\partial \xi_k}\,.
  \label{eq:dNi_vertex}
\end{equation}

For the edge midpoint $N_{ij} = 4\lambda_i\lambda_j$:
\begin{equation}
  \frac{\partial N_{ij}}{\partial \xi_k}
  = 4\Bigl(
    \lambda_j\,\frac{\partial \lambda_i}{\partial \xi_k}
  + \lambda_i\,\frac{\partial \lambda_j}{\partial \xi_k}
  \Bigr).
  \label{eq:dNij_edge}
\end{equation}


% ─────────────────────────────────────────────────────────────────────────
\section{Isoparametric Mapping}
\label{sec:isoparametric}

The isoparametric map from the reference simplex $\calS_D$ to a physical
element $\Omega^e \subset \R^d$ (with $d \ge D$) is
\begin{equation}
  \mathbf{x}(\bxi)
  = \sum_{I=1}^{n_{\text{nodes}}} N_I(\bxi)\,\mathbf{x}_I,
  \label{eq:isoparametric_map}
\end{equation}
where $\mathbf{x}_I \in \R^d$ are the physical nodal coordinates.  The
\emph{Jacobian matrix} of the mapping is
\begin{equation}
  J_{ij} = \frac{\partial x_i}{\partial \xi_j}
         = \sum_{I=1}^{n_{\text{nodes}}}
           \frac{\partial N_I}{\partial \xi_j}\, x_{I,i}\,,
  \qquad
  \mathbf{J} \in \R^{d \times D}.
  \label{eq:jacobian}
\end{equation}

When $d = D$ (volume elements), $\mathbf{J}$ is square and the differential
volume element is $\dd\Omega = |\det\mathbf{J}|\,\dd\bxi$.

When $d > D$ (surface elements, e.g.\ TRI6 boundary faces in 3D), the surface
measure is computed from the metric tensor:
\begin{equation}
  \dd S = \sqrt{\det(\mathbf{J}^T \mathbf{J})}\;\dd\bxi.
  \label{eq:surface_measure}
\end{equation}

\paragraph{Gradient transformation.}
The gradient of a scalar field in physical space is obtained from the
reference-space gradient via
\begin{equation}
  \frac{\partial N_I}{\partial x_j}
  = \sum_{k=1}^{D} \frac{\partial N_I}{\partial \xi_k}\,
    (J^{-1})_{kj}\,.
  \label{eq:grad_transform}
\end{equation}


% ─────────────────────────────────────────────────────────────────────────
\section{Subentity Topology: Faces of the Simplex}
\label{sec:faces}

A $D$-simplex has $D+1$ faces (codimension-$1$ sub-simplices).  The $f$-th face
($f = 0, \dots, D$) is the sub-simplex \emph{opposite} vertex $f$, spanned by
all vertices \emph{except} $\mathbf{v}_f$.

\subsection{Face Local Node Ordering (TET10)}

For a quadratic tetrahedron ($D=3$, TET10, 10~nodes), each of the 4~triangular
faces is a TRI6 sub-element.  The face-to-global node mapping implemented in
\texttt{SimplexCell<3,2>} is:

\begin{center}
\begin{tabular}{ccc}
  \toprule
  Face $f$ & Opposite vertex & Global nodes \\
  \midrule
  0 & $\mathbf{v}_0$ & (1, 2, 3, 5, 8, 7) \\
  1 & $\mathbf{v}_1$ & (0, 3, 2, 6, 8, 5) \\
  2 & $\mathbf{v}_2$ & (0, 1, 3, 4, 7, 6) \\
  3 & $\mathbf{v}_3$ & (0, 2, 1, 5, 4, 6) \\
  \bottomrule
\end{tabular}
\end{center}

The outward normal direction follows from the right-hand rule applied to the
first three (vertex) nodes of each face.  This topology is critical for
consistent surface traction application via Gauss quadrature on boundary faces.


% ─────────────────────────────────────────────────────────────────────────
\section{Gmsh Node Ordering for TET10}
\label{sec:gmsh_reorder}

Gmsh uses a different midpoint ordering than
\texttt{fall\_n}~\cite{Geuzaine2009}.  The Gmsh TET10 local ordering is:

\begin{center}
\begin{tabular}{ccccccccccc}
  \toprule
  Position & 0 & 1 & 2 & 3 & 4 & 5 & 6 & 7 & 8 & 9 \\
  \midrule
  Gmsh     & $v_0$ & $v_1$ & $v_2$ & $v_3$ & $m_{01}$ & $m_{12}$ & $m_{02}$ & $m_{03}$ & $m_{13}$ & $m_{23}$ \\
  \texttt{fall\_n} & $v_0$ & $v_1$ & $v_2$ & $v_3$ & $m_{01}$ & $m_{02}$ & $m_{03}$ & $m_{12}$ & $m_{23}$ & $m_{13}$ \\
  \bottomrule
\end{tabular}
\end{center}

The reordering applied in \texttt{GmshDomainBuilder.hh} maps Gmsh indices to
\texttt{fall\_n} indices via the permutation:
\begin{equation}
  (0, 1, 2, 3, 4, 6, 7, 5, 9, 8).
  \label{eq:gmsh_perm}
\end{equation}
This is implemented as:
\begin{lstlisting}
index_ordering(node_tags, 0,1,2,3,4,6,7,5,9,8).data()
\end{lstlisting}


% ─────────────────────────────────────────────────────────────────────────
\section{Implementation in \texttt{fall\_n}}
\label{sec:simplex_impl}

The simplex element hierarchy consists of three layers:

\begin{enumerate}[leftmargin=2em]
  \item \textbf{\texttt{SimplexCell<D, Order>}} (\texttt{SimplexCell.hh}) ---
        Purely geometric: reference coordinates, shape functions and their
        derivatives, face topology.  All methods are \texttt{consteval} or
        \texttt{constexpr}, enabling compile-time evaluation.

  \item \textbf{\texttt{SimplexElement<Dim, TopDim, Order>}}
        (\texttt{SimplexElement.hh}) --- Combines a \texttt{SimplexCell} with
        an integration strategy.  Template parameters:
        \begin{itemize}[nosep]
          \item \texttt{Dim}: ambient space dimension (e.g., 3 for 3D problems);
          \item \texttt{TopDim}: topological dimension ($D$);
          \item \texttt{Order}: Lagrange interpolation order.
        \end{itemize}

  \item \textbf{\texttt{ElementGeometry<ElemType, Integrator>}}
        (\texttt{ElementGeometry.hh}) --- Manages the Jacobian, physical
        gradients, and integration; parameterised by element type and
        quadrature rule.  The integrator is a duck-typed object exposing
        \texttt{num\_integration\_points}, \texttt{reference\_integration\_point(i)},
        and \texttt{weight(i)}.
\end{enumerate}

This layered design allows the quadrature rule to be swapped independently of
the element geometry, which was crucial for diagnosing and resolving the
negative-weight problem (see~\cref{ch:negative_weights,ch:quadrature}).

% =============================================================================
%  Chapter 2 — Numerical Quadrature on Simplices
% =============================================================================

\chapter{Numerical Quadrature on Simplices}
\label{ch:quadrature}

Numerical integration on simplicial domains is fundamentally different from
integration on tensor-product cells (quadrilaterals, hexahedra).  On
tensor-product cells, one constructs $n^D$-point rules by taking the Cartesian
product of $n$-point Gauss--Legendre rules in each direction.  On simplices, no
such natural product structure exists; special rules must be
constructed~\cite{Stroud1971,Keast1986}.

This chapter surveys the approaches available, documents the rules used in
\texttt{fall\_n}, and motivates the Stroud conical product as the preferred
arbitrary-order method.


% ─────────────────────────────────────────────────────────────────────────
\section{Problem Statement}
\label{sec:quad_problem}

We seek to approximate the integral
\begin{equation}
  I[f] = \int_{\calS_D} f(\bxi)\,\dd\bxi
  \approx \sum_{g=1}^{n_g} w_g\, f(\bxi_g),
  \label{eq:quad_approx}
\end{equation}
where $\calS_D$ is the reference $D$-simplex~\eqref{eq:ref_simplex}, with
volume $|\calS_D| = 1/D!$.  A rule has \emph{degree of exactness}~$p$ if
$I[f]$ is computed exactly for all polynomials $f$ of total degree $\le p$.

\paragraph{Minimum requirements for finite elements.}
For a displacement-based finite element with polynomial order $k$ and a
constitutive law that is polynomial of degree $m$ in the strains:
\begin{itemize}[nosep]
  \item Stiffness matrix: needs degree $\ge 2(k-1)m$ (often $2(k-1)$ for
        linear elasticity, where $m=1$).
  \item Mass matrix (or $L^2$ projection): needs degree $\ge 2k$.
  \item For TET10 ($k=2$, $D=3$): stiffness needs degree~$\ge 2$; mass/projection
        needs degree~$\ge 4$.
\end{itemize}


% ─────────────────────────────────────────────────────────────────────────
\section{Classical Symmetric Rules}
\label{sec:classical_rules}

\subsection{Centroid Rule (Degree 1)}

The single-point centroid rule:
\begin{equation}
  \bxi_1 = \Bigl(\frac{1}{D+1}, \dots, \frac{1}{D+1}\Bigr),
  \qquad w_1 = \frac{1}{D!}\,.
  \label{eq:centroid_rule}
\end{equation}
This integrates all linear polynomials exactly.  It is adequate for constant-
strain elements (TET4) when only the stiffness matrix is needed, but
insufficient for quadratic elements.

\subsection{Hammer--Marlowe--Stroud Rules (HMS)}
\label{sec:HMS}

The Hammer--Marlowe--Stroud~\cite{HammerMarloweStroud1956} family provides low-
order symmetric rules.  The 4-point rule for the tetrahedron, used as the
default for TET10 in \texttt{fall\_n}, has:

\begin{center}
\begin{tabular}{ccl}
  \toprule
  $n_g$ & Degree & Description \\
  \midrule
  1     & 1      & Centroid \\
  4     & 2      & Vertices of a smaller inscribed tetrahedron \\
  \bottomrule
\end{tabular}
\end{center}

The 4-point degree-2 rule uses barycentric coordinates:
\begin{equation}
  \lambda = \bigl(\alpha, \beta, \beta, \beta\bigr)
  \quad\text{and permutations thereof},
  \qquad
  \alpha = \frac{5 - \sqrt{5}}{20},\;
  \beta  = \frac{5 + 3\sqrt{5}}{20}.
\end{equation}
All four weights are equal, $w_g = \frac{1}{4!} = \frac{1}{24}$, and are
\textbf{strictly positive}.

\begin{remark}
In \texttt{fall\_n}, this 4-point HMS rule is the \emph{default} quadrature
for TET10 elements (\texttt{default\_simplex\_rule<3,2>()} in
\texttt{SimplexQuadrature.hh}).  It was chosen over the 5-point Keast degree-3
rule precisely because all weights are positive\,---\,see
\cref{ch:negative_weights} for the detailed analysis.
\end{remark}


\subsection{Keast Rules}
\label{sec:keast}

Keast~\cite{Keast1986} tabulated a series of tetrahedral rules up to degree~8.
The most commonly cited are:

\begin{center}
\begin{tabular}{cccl}
  \toprule
  $n_g$ & Degree & Min weight & Comment \\
  \midrule
  1     & 1    & $+$      & Centroid \\
  4     & 2    & $+$      & Same as HMS \\
  5     & 3    & $\mathbf{-}$ & $w_1 = -\tfrac{2}{15}$ at centroid \\
  11    & 4    & $+$      & All positive \\
  14    & 5    & $+$      & All positive \\
  24    & 6    & $\mathbf{-}$ & Contains negative weights \\
  \bottomrule
\end{tabular}
\end{center}

The 5-point degree-3 rule is particularly notable:
\begin{equation}
  \text{Keast-5:}\quad
  w_1 = -\frac{2}{15}\,|\calS_3| = -\frac{2}{15} \cdot \frac{1}{6}
      = -\frac{1}{45},\qquad
  w_{2,\dots,5} = \frac{3}{40}\,|\calS_3| = \frac{1}{80}.
  \label{eq:keast5_weights}
\end{equation}
The centroid weight is \emph{negative}.  While this does not affect the
accuracy of the quadrature rule (it still integrates degree-3 polynomials
exactly), it causes severe problems in lumped $L^2$ projection on quadratic
elements, as analysed in \cref{ch:negative_weights}.


\subsection{Grundmann--M\"oller Rules}

The Grundmann--M\"oller family~\cite{GrundmannMoller1978} generates rules of
arbitrary odd degree for any simplex dimension, using a combinatorial formula.
The rule of index $s$ has degree $2s+1$ with
\[
  n_g = \binom{D + s + 1}{D + 1}
\]
points.  Unfortunately, for $s \ge 1$ (degree $\ge 3$), some weights are
\emph{guaranteed to be negative}.  This makes Grundmann--M\"oller unsuitable
as a general-purpose rule when positive weights are required.


% ─────────────────────────────────────────────────────────────────────────
\section{Stroud Conical Product Rule}
\label{sec:conical_product}

The Stroud conical product~\cite{Stroud1971,Duffy1982} transforms the simplex
integral into a product of 1D integrals via a \emph{collapsed-coordinate}
(Duffy) transformation, enabling the use of Gauss--Jacobi quadrature in each
direction.


\subsection{Duffy Transform}
\label{sec:duffy}

The transformation from \emph{collapsed coordinates}
$(t_1, \dots, t_D) \in [0,1]^D$ to simplex coordinates
$(\xi_1, \dots, \xi_D) \in \calS_D$ is defined recursively:
\begin{equation}
  \boxed{
  \begin{aligned}
    \xi_1 &= t_1, \\
    \xi_k &= t_k \prod_{j=1}^{k-1}(1 - t_j), \qquad k = 2, \dots, D.
  \end{aligned}}
  \label{eq:duffy_transform}
\end{equation}

The Jacobian of this transformation is
\begin{equation}
  \left|\frac{\partial\bxi}{\partial\mathbf{t}}\right|
  = \prod_{k=1}^{D-1} (1 - t_k)^{D-k}.
  \label{eq:duffy_jacobian}
\end{equation}

Therefore, the integral over the simplex becomes
\begin{equation}
  \int_{\calS_D} f(\bxi)\,\dd\bxi
  = \int_0^1 \cdots \int_0^1
    f\bigl(\bxi(\mathbf{t})\bigr)
    \prod_{k=1}^{D-1} (1 - t_k)^{D-k}
    \;\dd t_1 \cdots \dd t_D.
  \label{eq:duffy_integral}
\end{equation}


\subsection{Product Quadrature with Gauss--Jacobi Nodes}

The key insight is that each factor $(1 - t_k)^{D-k}$ depends only on $t_k$,
so it can be \emph{absorbed into} the 1D quadrature rule for direction $k$.
Specifically, for direction $k$ ($0$-indexed), we need a 1D quadrature rule on
$[0,1]$ with weight function $(1-t)^{\beta_k}$, where
\begin{equation}
  \beta_k = D - 1 - k, \qquad k = 0, 1, \dots, D-1.
  \label{eq:beta_definition}
\end{equation}

\begin{itemize}[nosep]
  \item $k = 0$: $\beta = D-1$ (maximal weight for the ``lowest'' direction).
  \item $k = D-2$: $\beta = 1$ (linear weight).
  \item $k = D-1$: $\beta = 0$ (plain Gauss--Legendre on $[0,1]$).
\end{itemize}

The integral~\eqref{eq:duffy_integral} becomes
\begin{equation}
  \int_{\calS_D} f(\bxi)\,\dd\bxi
  = \prod_{k=0}^{D-1} \left(
    \int_0^1 g_k(t_k)\,(1 - t_k)^{\beta_k}\,\dd t_k
  \right),
  \label{eq:product_integral}
\end{equation}
where $g_k$ encodes the dependence of $f$ on $t_k$ after integrating all other
variables.  Applying $n$-point Gauss--Jacobi quadrature on $[0,1]$ with weight
$(1-t)^{\beta_k}$ to each direction gives a product rule with:
\begin{equation}
  \boxed{
  n_g = n^D \text{ points}, \qquad
  \text{degree of exactness} = 2n - 1.}
  \label{eq:conical_cost}
\end{equation}

\begin{proposition}[All weights positive]
\label{prop:positive_weights}
Since each 1D Gauss--Jacobi rule has strictly positive weights (the weight
function $(1-t)^{\beta}$ is positive on $(0,1)$ for $\beta > -1$, and all
Gauss--Jacobi weights are positive for any $\alpha, \beta > -1$), the product
rule inherits positivity:
\begin{equation}
  w_g = \prod_{k=0}^{D-1} w_{k,i_k} > 0 \qquad \forall\, g.
  \label{eq:positive_product_weight}
\end{equation}
This is the fundamental advantage of the conical product over symmetric
simplex rules.
\end{proposition}


\subsection{Cost Comparison}

\begin{center}
\begin{tabular}{lccl}
  \toprule
  Rule & Points & Degree & Weights \\
  \midrule
  HMS 4-pt (TET)          & 4    & 2 & All $+$ \\
  Keast 5-pt              & 5    & 3 & \textbf{1 negative} \\
  Keast 11-pt             & 11   & 4 & All $+$ \\
  Conical $n\!=\!2$, $D\!=\!3$ & 8   & 3 & All $+$ \\
  Conical $n\!=\!3$, $D\!=\!3$ & 27  & 5 & All $+$ \\
  Conical $n\!=\!4$, $D\!=\!3$ & 64  & 7 & All $+$ \\
  \bottomrule
\end{tabular}
\end{center}

The conical product uses more points than optimal symmetric rules for the same
degree, but it guarantees positivity and extends to \emph{any} polynomial
degree.  For the TET10 stiffness matrix (degree $\ge 2$), the 4-point HMS rule
is sufficient and more economical; the conical product is available as a
drop-in replacement when higher accuracy is needed.


% ─────────────────────────────────────────────────────────────────────────
\section{Gauss--Jacobi Quadrature via the Golub--Welsch Algorithm}
\label{sec:golub_welsch}

The 1D Gauss--Jacobi rules needed for the conical product are computed at
library initialisation time using the Golub--Welsch
algorithm~\cite{GolubWelsch1969}.

\subsection{Jacobi Polynomials}
\label{sec:jacobi_poly}

The Jacobi polynomials $P_n^{(\alpha,\beta)}(x)$ are orthogonal on $[-1,1]$
with respect to the weight function
\begin{equation}
  w(x) = (1 - x)^\alpha (1 + x)^\beta, \qquad \alpha, \beta > -1.
  \label{eq:jacobi_weight}
\end{equation}
They satisfy the three-term recurrence~\cite{Szego1975,Gautschi2004}:
\begin{equation}
  P_{n+1}^{(\alpha,\beta)}(x)
  = (A_n\, x + B_n)\, P_n^{(\alpha,\beta)}(x)
  - C_n\, P_{n-1}^{(\alpha,\beta)}(x),
  \label{eq:3term_recurrence}
\end{equation}
where the coefficients $A_n$, $B_n$, $C_n$ are explicit functions of $n$,
$\alpha$, $\beta$.

The \emph{monic} Jacobi polynomials $\hat{p}_n$ satisfy
\[
  \hat{p}_{n+1}(x) = (x - a_n)\,\hat{p}_n(x) - b_n\,\hat{p}_{n-1}(x),
\]
with recurrence coefficients~\cite{Gautschi2004}:
\begin{align}
  a_n &= \frac{\beta^2 - \alpha^2}{(2n+\alpha+\beta)(2n+\alpha+\beta+2)},
  \label{eq:an_coeff} \\
  b_n &= \frac{4n(n+\alpha)(n+\beta)(n+\alpha+\beta)}
              {(2n+\alpha+\beta)^2(2n+\alpha+\beta+1)(2n+\alpha+\beta-1)}.
  \label{eq:bn_coeff}
\end{align}


\subsection{The Golub--Welsch Algorithm}
\label{sec:gw_algorithm}

The Golub--Welsch algorithm~\cite{GolubWelsch1969} computes the $n$-point
Gauss quadrature rule from the eigendecomposition of the \emph{symmetric
tridiagonal Jacobi matrix} $\mathbf{T} \in \R^{n \times n}$:
\begin{equation}
  \mathbf{T} = \begin{pmatrix}
    a_0       & \sqrt{b_1} &            &        \\
    \sqrt{b_1} & a_1       & \sqrt{b_2} &        \\
                & \ddots    & \ddots     & \ddots \\
                &           & \sqrt{b_{n-1}} & a_{n-1}
  \end{pmatrix}.
  \label{eq:jacobi_matrix}
\end{equation}

\begin{itemize}[nosep]
  \item The \textbf{eigenvalues} $\lambda_1, \dots, \lambda_n$ of
        $\mathbf{T}$ are the \emph{quadrature nodes} (zeros of the degree-$n$
        monic Jacobi polynomial).
  \item The \textbf{weights} are $w_i = \mu_0\, v_{1,i}^2$, where
        $v_{1,i}$ is the first component of the $i$-th normalised eigenvector
        and $\mu_0$ is the zeroth moment of the weight function:
        \begin{equation}
          \mu_0 = \int_{-1}^{1} (1-x)^\alpha (1+x)^\beta\,\dd x
               = \frac{2^{\alpha+\beta+1}\,\Gamma(\alpha+1)\,\Gamma(\beta+1)}
                      {\Gamma(\alpha+\beta+2)}.
          \label{eq:mu0}
        \end{equation}
\end{itemize}

\begin{remark}[Why Golub--Welsch instead of Newton iteration]
An earlier implementation attempted to find the Jacobi polynomial zeros via a
Newton iteration with Tricomi-type initial guesses.  This approach failed for
the $\alpha = D-1$ cases needed by the conical product (e.g., $\alpha = 2$ for
$D = 3$, direction $k=0$): the initial guesses were too inaccurate, and the
iteration converged to \emph{duplicate roots} instead of distinct nodes.

The Golub--Welsch algorithm is unconditionally robust: it requires no initial
guesses and produces all nodes and weights simultaneously from a single
eigenvalue decomposition.  It is implemented using Eigen's
\texttt{SelfAdjointEigenSolver}, which is $O(n^2)$ for the tridiagonal case.
\end{remark}


\subsection{Mapping to $[0,1]$}

The Gauss--Jacobi rule on $[-1,1]$ with weight $(1-x)^\alpha$ is mapped to
$[0,1]$ via $x = 2t - 1$:
\begin{equation}
  \int_0^1 f(t)\,(1-t)^\beta\,\dd t
  = \frac{1}{2^{\beta+1}} \int_{-1}^{1} f\!\left(\frac{1+x}{2}\right)
    (1-x)^\beta\,\dd x.
  \label{eq:map_to_01}
\end{equation}
The nodes and weights transform as
\begin{equation}
  t_i = \frac{1 + x_i}{2}, \qquad
  w_i^{[0,1]} = \frac{w_i^{[-1,1]}}{2^{\beta+1}}.
  \label{eq:weight_map}
\end{equation}
This is implemented in \texttt{gauss\_jacobi\_01()} in
\texttt{StroudConicalProduct.hh}.


% ─────────────────────────────────────────────────────────────────────────
\section{Implementation: \texttt{StroudConicalProduct.hh}}
\label{sec:conical_impl}

The implementation resides in
\texttt{src/numerics/numerical\_integration/StroudConicalProduct.hh} and
consists of three layers:

\begin{enumerate}[leftmargin=2em]
  \item \textbf{\texttt{gauss\_jacobi(n, $\alpha$, $\beta$)}} ---
        Returns an $n$-point Gauss--Jacobi rule on $[-1,1]$ via Golub--Welsch.
        Uses Eigen's \texttt{SelfAdjointEigenSolver} on the tridiagonal Jacobi
        matrix.

  \item \textbf{\texttt{stroud\_conical\_product<D>(n\_per\_dir)}} ---
        Builds $D$ one-dimensional Gauss--Jacobi rules (with
        $\beta_k = D-1-k$), takes their tensor product in collapsed
        coordinates, and maps to simplex coordinates via the Duffy
        transform~\eqref{eq:duffy_transform}.

  \item \textbf{\texttt{ConicalProductIntegrator<TopDim, NPerDir>}} ---
        Compile-time-sized integrator class.  The rule is computed once at
        program startup (static \texttt{inline const} initialiser) and stored
        in a fixed-size \texttt{std::array}.  Exposes the same interface as
        all other integrators in \texttt{fall\_n}:
        \begin{itemize}[nosep]
          \item \texttt{static constexpr num\_integration\_points = NPerDir$^{\text{TopDim}}$}
          \item \texttt{reference\_integration\_point(i)} returns a $\texttt{std::span<const double>}$
          \item \texttt{weight(i)} returns $w_i$
          \item \texttt{operator()(f)} evaluates $\sum w_i f(\bxi_i)$
        \end{itemize}
\end{enumerate}


% ─────────────────────────────────────────────────────────────────────────
\section{Default Quadrature Choices in \texttt{fall\_n}}
\label{sec:default_rules}

\begin{center}
\begin{tabular}{llccl}
  \toprule
  Element & Rule & Points & Degree & Source \\
  \midrule
  TET4  & HMS 1-pt  & 1  & 1 & \texttt{default\_simplex\_rule<3,1>()} \\
  TET10 & HMS 4-pt  & 4  & 2 & \texttt{default\_simplex\_rule<3,2>()} \\
  TRI3  & HMS 1-pt  & 1  & 1 & \texttt{default\_simplex\_rule<2,1>()} \\
  TRI6  & HMS 3-pt  & 3  & 2 & \texttt{default\_simplex\_rule<2,2>()} \\
  HEX27 & Gauss $3^3$ & 27 & 5 & Tensor-product Gauss--Legendre \\
  \bottomrule
\end{tabular}
\end{center}

The conical product integrators (\texttt{ConicalProductIntegrator<3,2>},
\texttt{ConicalProductIntegrator<3,3>}, etc.) are available as drop-in
replacements when higher-degree or higher-accuracy integration is needed,
with the guarantee that all weights remain positive.

% =============================================================================
%  Chapter 3 — Negative Weights and the Checkerboard Artefact
% =============================================================================

\chapter{Negative Weights and the Checkerboard Artefact}
\label{ch:negative_weights}

This chapter documents the diagnosis and resolution of a persistent
\emph{checkerboard} (or ``diamond'') artefact observed in nodal stress fields
projected from Gauss points onto quadratic tetrahedral elements (TET10).  The
analysis traces the root cause to the interaction between negative quadrature
weights and shape function values, and prescribes two complementary fixes.


% ─────────────────────────────────────────────────────────────────────────
\section{Symptom: Element-Level Checkerboard Pattern}
\label{sec:symptom}

When visualising projected stress or strain fields in ParaView on TET10 meshes,
the following pathology was observed:
\begin{itemize}[nosep]
  \item At \emph{vertex} nodes, the projected field values had the
        \emph{opposite sign} of the expected physical values.
  \item At \emph{edge midpoint} nodes, the values were correct.
  \item The pattern repeated \emph{at every element}, forming a regular
        alternating (checkerboard) structure.
  \item On HEX27 meshes, the same analysis pipeline produced smooth, correct
        fields.
\end{itemize}


% ─────────────────────────────────────────────────────────────────────────
\section{Root Cause Analysis}
\label{sec:root_cause}

The projected nodal values were computed via \emph{lumped $L^2$ projection}:
\begin{equation}
  \tilde{\sigma}_I
  = \frac{\displaystyle\sum_{e \in \mathcal{E}_I}
          \sum_{g=1}^{n_g} N_I(\bxi_g)\, w_g\, |J_g|\, \sigma_g}
         {\displaystyle\sum_{e \in \mathcal{E}_I}
          \sum_{g=1}^{n_g} N_I(\bxi_g)\, w_g\, |J_g|},
  \label{eq:lumped_L2}
\end{equation}
where $\mathcal{E}_I$ is the set of elements sharing node $I$, $N_I$ is the
global shape function, and $\sigma_g$ is the Gauss-point value.

For the projection~\eqref{eq:lumped_L2} to produce meaningful results, the
denominator must be \textbf{strictly positive} for every node.  We will show
that this condition \emph{fails} at vertex nodes of TET10 when the Keast
5-point rule is used.


\subsection{The Keast 5-Point Rule Weights}

The Keast degree-3 rule for the tetrahedron has 5~points
(\cref{sec:keast}):
\begin{equation}
  \begin{aligned}
    &\text{Point 1 (centroid):}   & w_1 &= -\frac{2}{15}\,|\calS_3|
      = -\frac{1}{45} \approx -0.02222, \\
    &\text{Points 2--5 (vertices):} & w_g &= +\frac{3}{40}\,|\calS_3|
      = +\frac{1}{80} = 0.01250.
  \end{aligned}
  \label{eq:keast5_detail}
\end{equation}
The centroid weight $w_1$ is \textbf{negative}.


\subsection{TET10 Vertex Shape Functions at Quadrature Points}

The vertex shape function for node $i$ is (cf.\
\cref{eq:shape_vertex_quad}):
\[
  N_i(\bxi) = \lambda_i\,(2\lambda_i - 1).
\]
At the 5 Keast quadrature points:

\paragraph{At the centroid} ($\lambda_i = 1/4$ for all $i$):
\[
  N_i = \tfrac{1}{4}(2 \cdot \tfrac{1}{4} - 1) = \tfrac{1}{4}(-\tfrac{1}{2})
      = -\tfrac{1}{8}.
\]

\paragraph{At the vertex-orbit points:} Each vertex-orbit point has one
``dominant'' barycentric coordinate near $0.5854$ and three ``small'' ones near
$0.1382$.

For the vertex $i$ that is \emph{dominant} at point $g$:
\[
  N_i \approx 0.5854 \times (2 \times 0.5854 - 1) = 0.5854 \times 0.1708
      \approx 0.1000.
\]

For a vertex $i$ that is \emph{not dominant}:
\[
  N_i \approx 0.1382 \times (2 \times 0.1382 - 1) = 0.1382 \times (-0.7236)
      \approx -0.1000.
\]


\subsection{The Denominator Catastrophe}

Consider the denominator of~\eqref{eq:lumped_L2} for a \emph{vertex} node $I$,
restricted to a single element:
\begin{equation}
  D_I = \sum_{g=1}^{5} N_I(\bxi_g)\, w_g\, |J_g|.
  \label{eq:denominator_element}
\end{equation}

Assuming approximately constant $|J_g| \approx |J|$ within the element:
\begin{align}
  D_I &\approx |J| \bigl[
    \underbrace{N_I(\bxi_1)\,w_1}_{\text{centroid: } (-1/8)(-1/45) > 0}
    +\, N_I(\bxi_2)\,w_2
    + \cdots
    + N_I(\bxi_5)\,w_5
  \bigr].
  \label{eq:D_I_expansion}
\end{align}

The centroid contribution is \emph{positive} (negative times negative), but the
vertex-orbit contributions produce a mix of positive and negative terms.  The
crucial issue is that at the three vertex-orbit points where $N_I < 0$
(because $\lambda_I < 1/2$), the \emph{positive} weight $w_g > 0$ creates a
negative contribution, while at the one vertex-orbit point where $N_I > 0$
(the ``own'' vertex direction), the positive weight creates a positive
contribution.

When these are summed, the negative contributions can \emph{overwhelm} the
positive ones, making $D_I < 0$ or $D_I \approx 0$.


\subsection{Consequence: Sign Flip or Division by Zero}

\begin{equation}
  D_I \le 0 \implies \tilde{\sigma}_I
  = \frac{\text{(numerator)}}{D_I}
  \begin{cases}
    \text{has wrong sign}  & \text{if } D_I < 0, \\
    \text{diverges}        & \text{if } D_I = 0.
  \end{cases}
  \label{eq:sign_flip}
\end{equation}

In practice, the code included a guard $|D_I| < \epsilon \implies
\tilde{\sigma}_I = 0$.  This zeroed out vertex values while midpoint values
(which have $N_{ij} = 4\lambda_i\lambda_j \ge 0$ everywhere) remained correct.
The result was the characteristic checkerboard pattern:

\begin{center}
\begin{tabular}{lcc}
  \toprule
  Node type & Denominator $D_I$ & Projected value \\
  \midrule
  Vertex   & $\le 0$ or $\approx 0$ & $0$ or wrong sign \\
  Midpoint & $> 0$                   & Correct \\
  \bottomrule
\end{tabular}
\end{center}


% ─────────────────────────────────────────────────────────────────────────
\section{Why HEX27 Is Unaffected}
\label{sec:hex27_safe}

HEX27 uses the tensor-product $3 \times 3 \times 3$ Gauss--Legendre rule
(27~points), which has \textbf{all positive weights}.  Furthermore, the
Lagrangian shape functions on the reference cube $[-1,1]^3$ are products of 1D
Lagrange polynomials, and while they can go negative, they do so in a balanced
way that keeps the lumped denominator positive.  The combination of all-positive
weights and the tensor-product structure prevents the denominator catastrophe.


% ─────────────────────────────────────────────────────────────────────────
\section{Resolution: Two Complementary Fixes}
\label{sec:resolution}

Two independent corrections were applied:

\subsection{Fix 1: Switch to an All-Positive-Weight Quadrature Rule}

The root cause is the negative weight $w_1 = -2/15\,|\calS_3|$ in the Keast
5-point rule.  Replacing it with the \textbf{HMS 4-point degree-2 rule}, which
has all positive weights (\cref{sec:HMS}), eliminates the negative-weight
pathway entirely.

In \texttt{SimplexQuadrature.hh}, the default for TET10 was changed:
\begin{lstlisting}
template<>
inline constexpr auto default_simplex_rule<3, 2>() {
    // Was: make_simplex_rule_3d_5()  [Keast 5-pt, w1 = -2/15]
    // Now: make_simplex_rule_3d_4()  [HMS 4-pt, all w > 0]
    return make_simplex_rule_3d_4();
}
\end{lstlisting}

\begin{remark}
The HMS 4-point rule has degree~2, one degree lower than the Keast 5-point
rule (degree~3).  For the TET10 stiffness matrix, degree~$2$ is the minimum
requirement ($2(k-1) = 2$ for $k=2$), so accuracy is maintained for linear
elasticity.  For the mass matrix or high-order nonlinear constitutive models,
the conical product (\texttt{ConicalProductIntegrator<3,2>}, 8~points, degree~3)
is available as a higher-degree drop-in replacement with all positive weights.
\end{remark}

\subsection{Fix 2: Guard Against Negative Denominators}

As a defensive measure, the lumped $L^2$ projection code was modified to use
the absolute value of the denominator:
\begin{lstlisting}
denominator = std::abs(raw_denominator);
if (denominator < epsilon) denominator = epsilon;
\end{lstlisting}

This ensures that even if a future quadrature rule with negative weights is
inadvertently used, the projection produces a \emph{smoothed} (not
sign-flipped) field.  The \texttt{std::abs} guard does not affect the result
when all weights are positive, since the denominator is then guaranteed
positive.

\subsection{Fix 3 (Recommended): SPR Volume-Weighted Averaging}

The most robust solution is to bypass lumped $L^2$ projection entirely and use
\textbf{volume-weighted nodal averaging}\,---\,a simplified form of the
Superconvergent Patch Recovery (SPR) technique~\cite{ZienkiewiczZhu1992a}.
This is discussed in detail in \cref{ch:projection}.


% ─────────────────────────────────────────────────────────────────────────
\section{Summary}

The checkerboard artefact results from the interaction of three factors:

\begin{enumerate}[nosep]
  \item \textbf{Negative quadrature weight} in the Keast 5-point rule
        ($w_1 = -2/15\,|\calS_3|$).
  \item \textbf{Negative vertex shape functions} of TET10 at interior
        quadrature points ($N_i = \lambda_i(2\lambda_i - 1) < 0$ when
        $\lambda_i < 1/2$).
  \item \textbf{Lumped $L^2$ projection} with a denominator that sums signed
        products $N_I \cdot w_g$, allowing cancellation to zero or sign
        reversal at vertex nodes.
\end{enumerate}

Any one of these factors being absent would prevent the artefact.  The chosen
resolution addresses all three:
\begin{itemize}[nosep]
  \item Factor~1: eliminated by switching to HMS 4-point (all $w_g > 0$);
  \item Factor~3: eliminated by switching to SPR averaging (no denominator
        subtraction possible).
  \item Factor~2: is intrinsic to quadratic elements and cannot be changed, but
        it is rendered harmless by the other two fixes.
\end{itemize}

% =============================================================================
%  Chapter 4 — Gauss-Point to Nodal Projection Methods
% =============================================================================

\chapter{Gauss-Point to Nodal Projection}
\label{ch:projection}

In displacement-based finite elements, the primary unknowns (displacements)
are nodal.  Derived quantities\,---\,strains, stresses, internal
variables\,---\,are computed at \emph{Gauss (integration) points} inside each
element.  For post-processing and visualisation, these quantities must be
\emph{projected} or \emph{averaged} to the mesh nodes so that they can be
rendered as continuous fields in VTK/ParaView.

This chapter compares the three projection methods implemented in
\texttt{fall\_n} and explains the current adaptive default: lumped $L^2$ when
the local nodal lumping is strictly positive, and volume-weighted averaging
otherwise.


% ─────────────────────────────────────────────────────────────────────────
\section{Lumped \texorpdfstring{$L^2$}{L2} Projection}
\label{sec:l2_projection}

\subsection{Consistent \texorpdfstring{$L^2$}{L2} Projection}

The mathematically rigorous approach minimises the $L^2$ norm of the error
between the discrete Gauss-point field and a continuous nodal field.  Let
$\sigma^h$ be the Gauss-point field and $\tilde{\sigma}^h = \sum_I
\tilde{\sigma}_I N_I$ the nodal projection.  The consistent $L^2$ projection
solves:
\begin{equation}
  \mathbf{M}\,\tilde{\bsigma} = \mathbf{f},
  \label{eq:consistent_L2}
\end{equation}
where $\mathbf{M}$ is the global mass matrix ($M_{IJ} = \int_\Omega N_I N_J
\,\dd\Omega$) and $\mathbf{f}$ is the right-hand side ($f_I = \int_\Omega N_I
\sigma^h \,\dd\Omega$).

This requires solving a global linear system, which is expensive and may
introduce oscillations at element boundaries (Gibbs-like effects) for
discontinuous fields.

\subsection{Row-Sum Lumping}
\label{sec:lumping}

The standard shortcut is to \emph{lump} the mass matrix by replacing it with
its row sums on the diagonal:
\begin{equation}
  \tilde{M}_{II} = \sum_J M_{IJ}
  = \int_\Omega N_I \sum_J N_J \,\dd\Omega
  = \int_\Omega N_I\,\dd\Omega,
  \label{eq:lumped_mass}
\end{equation}
using the partition of unity $\sum_J N_J = 1$.  The lumped projection then
becomes a local (node-by-node) average:
\begin{equation}
  \tilde{\sigma}_I
  = \frac{\displaystyle\int_\Omega N_I\, \sigma^h \,\dd\Omega}
         {\displaystyle\int_\Omega N_I \,\dd\Omega}
  \approx
  \frac{\displaystyle\sum_{e \in \mathcal{E}_I}
        \sum_{g=1}^{n_g} N_I(\bxi_g)\, w_g\, |J_g|\, \sigma_g}
       {\displaystyle\sum_{e \in \mathcal{E}_I}
        \sum_{g=1}^{n_g} N_I(\bxi_g)\, w_g\, |J_g|}.
  \label{eq:lumped_L2_final}
\end{equation}

The denominator is the ``nodal volume'' contributed by the quadrature rule.
For it to be meaningful, we need
\begin{equation}
  D_I = \sum_{e \in \mathcal{E}_I}
        \sum_g N_I(\bxi_g)\, w_g\, |J_g| > 0
        \qquad \forall\, I.
  \label{eq:denominator_positive}
\end{equation}

As demonstrated in \cref{ch:negative_weights} and refined in
\cref{ch:adaptive_projection}, this condition can fail in two distinct ways:
\begin{itemize}[nosep]
  \item \emph{catastrophically}, when negative quadrature weights amplify the
        signed support of higher-order simplex basis functions;
  \item \emph{structurally}, when the exact row-sum lumping
        $\int_\Omega N_I\,\dd\Omega$ is itself non-positive, as happens for the
        vertex functions of TET10 even under all-positive quadrature.
\end{itemize}
In the second case the formula is still well-defined, but it is no longer a
nodal average.  It becomes a signed extrapolation operator, which is precisely
what produces the spurious oscillations seen in TET10 contour plots.

\paragraph{Implementation.}  The lumped $L^2$ projection is implemented in
\texttt{l2\_project\_to\_nodes()} in \texttt{VTKModelExporter.hh}.  The legacy
\texttt{std::abs} guard on the denominator remains in the explicit
\texttt{LumpedL2} path as a defensive measure against future negative-weight
rules:
\begin{lstlisting}
double denom = std::abs(node_lump[I]);
if (denom < 1e-14) denom = 1e-14;
node_val[I] /= denom;
\end{lstlisting}
but this is no longer relied on as the production fix for quadratic simplex
elements.  The production path is now the adaptive resolver documented in
\cref{ch:adaptive_projection}.


% ─────────────────────────────────────────────────────────────────────────
\section{Volume-Weighted Averaging (SPR-Type)}
\label{sec:spr_averaging}

\subsection{Motivation}

The Superconvergent Patch Recovery (SPR) method of Zienkiewicz and
Zhu~\cite{ZienkiewiczZhu1992a,ZienkiewiczZhu1992b} fits a polynomial to the
Gauss-point values in a \emph{patch} of elements surrounding each node and
evaluates the polynomial at the node.  In its full form, SPR requires solving
local least-squares problems.

For production post-processing where visual smoothness (rather than
superconvergent error estimation) is the primary goal, a simpler and even more
robust approach suffices: \textbf{volume-weighted element averaging}.

\subsection{Algorithm}
\label{sec:spr_algorithm}

For each node $I$, compute the nodal field as a weighted average of the
element centroid values:
\begin{equation}
  \boxed{
  \tilde{\sigma}_I
  = \frac{\displaystyle\sum_{e \in \mathcal{E}_I} V_e\, \bar{\sigma}_e}
         {\displaystyle\sum_{e \in \mathcal{E}_I} V_e},
  }
  \label{eq:spr_average}
\end{equation}
where:
\begin{itemize}[nosep]
  \item $\mathcal{E}_I$ is the set of elements sharing node $I$;
  \item $V_e$ is the volume of element $e$ (computed via quadrature:
        $V_e = \sum_g w_g |J_g|$);
  \item $\bar{\sigma}_e$ is the element-average value, computed as the
        weighted mean of Gauss-point values:
        \begin{equation}
          \bar{\sigma}_e
          = \frac{\sum_g w_g\, |J_g|\, \sigma_g}{\sum_g w_g\, |J_g|}.
          \label{eq:element_average}
        \end{equation}
\end{itemize}

\subsection{Robustness Guarantee}

\begin{proposition}[Unconditional positivity]
All weights in the averaging formula~\eqref{eq:spr_average} are element
volumes $V_e > 0$, which are strictly positive for any non-degenerate element
regardless of the quadrature rule used for stiffness integration.  Even if the
quadrature rule contains negative weights, the element volume $V_e = \sum_g
w_g |J_g|$ remains positive (it is the exact volume when the quadrature degree
is sufficient to integrate polynomials of the Jacobian's degree).

Therefore, SPR volume-weighted averaging is \textbf{unconditionally robust}:
the denominator is always positive, and no sign flip or division by zero can
occur at any node, regardless of element type or quadrature rule.
\end{proposition}

\subsection{Comparison with Lumped \texorpdfstring{$L^2$}{L2}}

\begin{center}
\renewcommand{\arraystretch}{1.2}
\begin{tabular}{p{4cm}cc}
  \toprule
  Property & Lumped $L^2$ & Volume-weighted avg. \\
  \midrule
  Denominator sign    & Can be $\le 0$ & Always $> 0$ \\
  Sensitivity to $w_g < 0$ & High (catastrophic) & None \\
  Accuracy order      & Optimal ($O(h^{k+1})$) & $O(h^k)$ \\
  Superconvergent?    & No  & No \\
  Computational cost  & $O(n_e \cdot n_g)$ & $O(n_e \cdot n_g)$ \\
  Requires global solve? & No & No \\
  \bottomrule
\end{tabular}
\end{center}

For visualisation purposes, the $O(h^k)$ vs.\ $O(h^{k+1})$ accuracy
difference is negligible.  The robustness advantage of volume-weighted
averaging is decisive.


% ─────────────────────────────────────────────────────────────────────────
\section{Implementation: \texttt{spr\_average\_to\_nodes()}}
\label{sec:spr_impl}

The volume-weighted averaging is implemented in two methods in
\texttt{VTKModelExporter.hh}:

\begin{description}
  \item[\texttt{spr\_average\_to\_nodes()}]
    Computes a single scalar field $\sigma$:
    \begin{enumerate}[nosep]
      \item For each element $e$, compute $V_e$ and $\bar{\sigma}_e$ from
            Gauss-point data.
      \item For each node $I$ of element $e$, accumulate:
            \begin{align}
              \text{numerator}[I]   &\mathrel{+}= V_e \cdot \bar{\sigma}_e, \\
              \text{denominator}[I] &\mathrel{+}= V_e.
            \end{align}
      \item After all elements: $\tilde{\sigma}_I = \text{numerator}[I] \,/\,
            \text{denominator}[I]$.
    \end{enumerate}

  \item[\texttt{compute\_material\_fields\_spr()}]
    Applies \texttt{spr\_average\_to\_nodes()} to all material fields: strain
    tensor (6 components), stress tensor (6 components), plastic strain tensor
    (6 components), and equivalent plastic strain (1 scalar).

  \item[\texttt{compute\_material\_fields()}]
    Resolves the projection strategy automatically.  It uses lumped $L^2$
    only if every element in the model has strictly positive local lumped nodal
    weights; otherwise it dispatches to volume-weighted averaging.
\end{description}


% ─────────────────────────────────────────────────────────────────────────
\section{Verification}
\label{sec:projection_verification}

The adaptive projection policy was verified on the cantilever benchmark
($10 \times 0.4 \times 0.8$ box, $E = 200$, $\nu = 0.3$):
\begin{itemize}[nosep]
  \item \textbf{HEX27 mesh} (20 elements, 315 nodes): the local lumped nodal
        weights are strictly positive, so the adaptive policy resolves to
        lumped $L^2$.  The projected fields remain smooth.
  \item \textbf{TET10 mesh} ($\sim 1329$ elements, 2840 nodes):
    \begin{itemize}[nosep]
      \item Lumped $L^2$ with Keast-5: \textbf{checkerboard at vertices}.
      \item Lumped $L^2$ with HMS-4 or Stroud conical product: still not a
            true nodal average, because the TET10 vertex lump weights remain
            negative.
      \item Adaptive mode therefore resolves to SPR volume-weighted
            averaging, which yields smooth, correct fields.
    \end{itemize}
  \item The main showcase (\texttt{main.cpp}) now uses
        \texttt{compute\_material\_fields()} with adaptive dispatch for all
        14~cases (7~formulations $\times$ 2~element types).
\end{itemize}

% =============================================================================
%  Chapter 5 — PETSc 3.24.5 Patch: section.c Indexing Bug
% =============================================================================

\chapter{PETSc Patch: \texttt{section.c} Indexing Bug}
\label{ch:petsc}

This chapter documents a bug in PETSc~3.24.5~\cite{PETSc2024} that affects
higher-order finite elements whose DMPlex point numbering does not start at
zero.  The bug was discovered during the development of simplex element support
in \texttt{fall\_n} and confirmed to be an upstream issue.


% ─────────────────────────────────────────────────────────────────────────
\section{Context: DMPlex and PetscSection}
\label{sec:dmplex_context}

PETSc's \texttt{DMPlex} represents an unstructured mesh as a directed acyclic
graph (DAG) of \emph{points}: vertices, edges, faces, and cells, numbered from
$0$ to $N-1$ (though the numbering of each \emph{stratum} may start at a
non-zero offset).

A \texttt{PetscSection} associates degrees of freedom (DOFs) with mesh points.
It uses an \emph{atlas} (\texttt{s->atlasDof[]}) indexed relative to the
chart $[\texttt{pStart}, \texttt{pEnd})$, and a \emph{constraint} mechanism
to impose Dirichlet boundary conditions by specifying which DOFs at a point
are constrained.

The function \texttt{PetscSectionSetConstraintIndices()} sets the constrained
DOF indices at a given point.  Its implementation reads:
\begin{lstlisting}[language=C]
PetscCall(PetscSectionCheckConstraints_Private(s));
/* Bug: should subtract pStart to get the atlas offset */
if (s->atlasDof[point]) {              // <-- HERE
    PetscInt cdof;
    PetscCall(PetscSectionGetConstraintDof(s, point, &cdof));
    // ...
}
\end{lstlisting}


% ─────────────────────────────────────────────────────────────────────────
\section{The Bug}
\label{sec:the_bug}

In \texttt{src/sys/utils/section.c}, at approximately line~2997, the code reads:
\begin{lstlisting}[language=C,escapechar=@]
if (s->atlasDof[@\textcolor{red}{point}@]) {
\end{lstlisting}
but the atlas array is indexed relative to \texttt{pStart}, so the correct
access is:
\begin{lstlisting}[language=C,escapechar=@]
if (s->atlasDof[@\textcolor{green!50!black}{point - s->pStart}@]) {
\end{lstlisting}

When \texttt{pStart = 0} (as happens for vertex stratum when vertices are
numbered from zero), the two expressions are identical and the bug is
invisible.  However, when a mesh's vertex stratum starts at a non-zero offset
--- which occurs in \texttt{fall\_n} because Gmsh-imported meshes place cells
before vertices in the DAG, giving the vertex stratum
$\texttt{pStart} = 8$~(for 8~hex cells) or larger --- the raw \texttt{point}
index overflows the atlas array, reading uninitialised memory.


% ─────────────────────────────────────────────────────────────────────────
\section{Manifestation in \texttt{fall\_n}}
\label{sec:manifestation}

The bug manifests during the application of Dirichlet boundary conditions
via \texttt{PetscSectionSetConstraintIndices}.  Specifically:
\begin{itemize}[nosep]
  \item For HEX27 meshes with 20 cells, the vertex stratum starts at
        \texttt{pStart} $\ge$ 20.
  \item When setting constraints on vertex DOFs, the index
        \texttt{s->atlasDof[point]} reads from memory at offset 20+
        instead of offset 0+.
  \item The read value is unpredictable: it may be zero (causing the
        constraints to be silently skipped) or non-zero (proceeding correctly
        by accident).
  \item The symptom was a segmentation fault or incorrect constraint
        application on some meshes but not others, depending on memory layout.
\end{itemize}


% ─────────────────────────────────────────────────────────────────────────
\section{The Fix}
\label{sec:the_fix}

The fix is a one-line change in
\texttt{src/sys/utils/section.c}:

\begin{lstlisting}[language=C]
// Before (buggy):
if (s->atlasDof[point]) {

// After (fixed):
if (s->atlasDof[point - s->pStart]) {
\end{lstlisting}

This was applied to the local PETSc installation at:
\begin{lstlisting}[basicstyle=\ttfamily\footnotesize,breaklines=true]
/home/sechavarriam/MyLibs/petsc-3.24.5/src/sys/utils/section.c
\end{lstlisting}

After applying the patch, PETSc was rebuilt and all \texttt{fall\_n} constraint
tests passed on both HEX27 and TET10 meshes.


% ─────────────────────────────────────────────────────────────────────────
\section{Verification}
\label{sec:petsc_verification}

The fix was verified by:
\begin{enumerate}[nosep]
  \item Running the debug output in \texttt{fall\_n}'s boundary condition
        application code, which prints:
\begin{lstlisting}[basicstyle=\ttfamily\footnotesize]
[SIEVE-DBG] section_ptr=0x... chart=[9152,23629)
[SIEVE-DBG] All 14477 section DOFs OK
[BC-DBG] Constraints count: 145
[BC-DBG] plex_id=9152 section_dof=3 constraint_size=3
  ...
\end{lstlisting}
        confirming that all 145~Dirichlet-constrained vertices are correctly
        handled with \texttt{pStart} = 9152.

  \item Comparing the Case~1 (SmallStrain + Linear Elastic) displacement field
        on HEX27 against the expected Timoshenko beam solution: the
        tip displacement matches the analytical result to within the expected
        discretisation error.

  \item Running all 7~nonlinear cases on both HEX27 and TET10 meshes without
        segfaults or constraint errors.
\end{enumerate}


% ─────────────────────────────────────────────────────────────────────────
\section{Upstream Status}

As of PETSc~3.24.5, this bug is present in the release source.  It has been
reported to the PETSc developers.  Since the bug only affects codes that:
\begin{itemize}[nosep]
  \item use \texttt{PetscSectionSetConstraintIndices} directly (rather than
        through higher-level \texttt{DMPlexSetBoundaryConditions}), and
  \item have a \texttt{PetscSection} chart with \texttt{pStart} $> 0$,
\end{itemize}
it may not have been triggered by the PETSc test suite, which may use meshes
with vertex-first ordering (\texttt{pStart} = 0 for the vertex stratum).

% =============================================================================
%  Chapter 6 — Implementation Summary and Design Decisions
% =============================================================================

\chapter{Implementation Summary and Design Decisions}
\label{ch:implementation}

This chapter collects the key design decisions made during the implementation
of simplex elements in \texttt{fall\_n}, with forward references to the
relevant source files and rationale grounded in the analysis of the preceding
chapters.


% ─────────────────────────────────────────────────────────────────────────
\section{Architecture Overview}
\label{sec:architecture}

\texttt{fall\_n} is a C++23 nonlinear finite element library.  Its element
hierarchy is fully \emph{compile-time polymorphic}: element types, quadrature
rules, constitutive models, and kinematic formulations are all resolved through
templates without virtual dispatch.

\begin{center}
\begin{tabular}{lll}
  \toprule
  Layer & Header & Responsibility \\
  \midrule
  Geometry        & \texttt{SimplexCell.hh}          & Reference coords, basis, topology \\
  Element         & \texttt{SimplexElement.hh}        & Iso-parametric map + integrator \\
  Integration     & \texttt{SimplexQuadrature.hh}     & Classical symmetric rules \\
                  & \texttt{StroudConicalProduct.hh}  & Arbitrary-order conical product \\
  Mesh I/O        & \texttt{GmshDomainBuilder.hh}     & Gmsh $\to$ \texttt{Domain} \\
  Model           & \texttt{Model.hh}                 & DOF management, BC application \\
  Post-processing & \texttt{VTKModelExporter.hh}      & Gauss $\to$ node projection, VTU \\
  \bottomrule
\end{tabular}
\end{center}


% ─────────────────────────────────────────────────────────────────────────
\section{Design Decisions}
\label{sec:decisions}


\subsection{Decision 1: HMS 4-Point as Default TET10 Quadrature}
\label{sec:dec_hms4}

\begin{description}[leftmargin=0pt,style=nextline]
  \item[Context] TET10 elements require at least degree-2 quadrature for the
    stiffness matrix.  The Keast 5-point rule (degree~3) was the initial
    default, following the common ``one degree above minimum'' heuristic.

  \item[Problem] The Keast 5-point rule has $w_1 = -2/15\,|\calS_3|$ at the
    centroid, causing the lumped $L^2$ projection to produce checkerboard
    artefacts (\cref{ch:negative_weights}).

  \item[Decision] Switch to HMS 4-point (degree~2) as the default.  Degree~2
    is the \emph{exact minimum} for the TET10 stiffness matrix; it provides
    no excess integration accuracy but avoids the negative-weight problem
    entirely.

  \item[Alternative considered] Conical product $n=2$ (8~points, degree~3,
    all positive).  This was rejected as the \emph{default} because it uses
    twice as many points as HMS 4, but it is available as
    \texttt{ConicalProductIntegrator<3,2>} for use cases requiring degree~3.

  \item[Source] \texttt{SimplexQuadrature.hh}, function
    \texttt{default\_simplex\_rule<3,2>()}.
\end{description}


\subsection{Decision 2: Golub--Welsch for Gauss--Jacobi Quadrature}
\label{sec:dec_gw}

\begin{description}[leftmargin=0pt,style=nextline]
  \item[Context] The Stroud conical product in 3D requires Gauss--Jacobi
    quadrature with parameters up to $\alpha = D-1 = 2$, $\beta = 0$.

  \item[Problem] An initial implementation used Newton iteration on the
    Jacobi polynomial (with Tricomi~/ asymptotic initial guesses).  For
    $\alpha = 2$ and $n \ge 3$, the iteration converged to duplicate roots
    because the initial guesses were insufficiently accurate for the
    distorted zero distribution of high-$\alpha$ Jacobi polynomials.

  \item[Decision] Replace Newton iteration entirely with the Golub--Welsch
    eigenvalue algorithm (\cref{sec:golub_welsch}).  This algorithm is
    unconditionally robust for all valid $\alpha, \beta > -1$ and any~$n$,
    since it requires no initial guesses.

  \item[Implementation] Uses Eigen's \texttt{SelfAdjointEigenSolver} on the
    $n \times n$ symmetric tridiagonal Jacobi matrix.  Cost is $O(n^2)$;
    since $n \le 10$ in practice, this is negligible.

  \item[Source] \texttt{StroudConicalProduct.hh}, function
    \texttt{gauss\_jacobi()}.
\end{description}


\subsection{Decision 3: Adaptive Projection as Default}
\label{sec:dec_spr}

\begin{description}[leftmargin=0pt,style=nextline]
  \item[Context] Gauss-point fields must be projected to nodes for VTU export.

  \item[Problem] Lumped $L^2$ projection is not robust on higher-order simplex
    elements.  The issue is deeper than negative quadrature weights:
    quadratic simplex vertex functions have signed support, so the row-sum
    lumped nodal measures can become non-positive even when \emph{all}
    quadrature weights are positive (\cref{ch:adaptive_projection}).

  \item[Decision] Make \texttt{compute\_material\_fields()} adaptive.  The
    exporter first inspects the local lumped nodal weights induced by the
    current element geometry and quadrature rule.  If every local weight is
    strictly positive, it uses lumped $L^2$ projection; otherwise it falls
    back to polynomial patch recovery
    (\texttt{patch\_recover\_to\_nodes()}).  The old volume-weighted averaging
    path is kept only as an explicit low-order fallback.

  \item[Trade-off] The adaptive pre-scan adds one post-processing pass over
    element quadrature points, but this cost is negligible compared with a
    solve and preserves the smoother lumped $L^2$ path for element families
    such as HEX27 where it is mathematically sound.  The patch-recovery path
    adds local least-squares work in post-processing only; it does not touch
    the FE assembly hot path.

  \item[Source] \texttt{VTKModelExporter.hh}, methods
    \texttt{lumped\_projection\_weights\_into()},
    \texttt{compute\_material\_fields()}, and
    \texttt{patch\_recover\_to\_nodes()}.
\end{description}


\subsection{Decision 4: Gmsh TET10 Node Reordering}
\label{sec:dec_reorder}

\begin{description}[leftmargin=0pt,style=nextline]
  \item[Context] Gmsh's local node ordering for 10-node tetrahedra differs
    from \texttt{fall\_n}'s ordering at the edge midpoints.

  \item[Problem] Without reordering, the Jacobian mapping is incorrect,
    producing negative determinants and garbage results.

  \item[Decision] Apply the permutation $(0,1,2,3,4,6,7,5,9,8)$ to map Gmsh
    indices to \texttt{fall\_n} indices.  This is done at mesh import time in
    \texttt{GmshDomainBuilder.hh}.

  \item[Verification] The simplex geometry tests
    (\texttt{test\_element\_geometry.cpp}) verify positive Jacobian
    determinants and correct isoparametric mapping for both TET4 and TET10
    elements imported from Gmsh.

  \item[Source] \texttt{GmshDomainBuilder.hh}, case 11 (TET10).
\end{description}


\subsection{Decision 5: PETSc Section Patch}
\label{sec:dec_petsc}

\begin{description}[leftmargin=0pt,style=nextline]
  \item[Context] \texttt{fall\_n} uses PETSc's DMPlex for distributed mesh
    management and \texttt{PetscSection} for DOF layout.

  \item[Problem] A bug in \texttt{PetscSectionSetConstraintIndices()} reads
    the atlas array with an incorrect index when \texttt{pStart} $> 0$
    (\cref{ch:petsc}).

  \item[Decision] Apply a local one-line patch:
    \texttt{s->atlasDof[point]} $\to$
    \texttt{s->atlasDof[point - s->pStart]}.

  \item[Impact] Without this patch, Dirichlet boundary conditions are silently
    incorrect on meshes where the vertex stratum does not start at point~0.

  \item[Source] PETSc~3.24.5, \texttt{src/sys/utils/section.c}, line~$\sim$2997.
\end{description}


\subsection{Decision 6: Hex27 VTK Face-Center Ordering}
\label{sec:dec_vtk_hex27}

\begin{description}[leftmargin=0pt,style=nextline]
  \item[Context] \texttt{fall\_n} stores the 27-node tensor-product HEX27 in
    lexicographic tensor order, while ParaView reads the exported VTU using
    \texttt{VTK\_TRIQUADRATIC\_HEXAHEDRON}.

  \item[Problem] The VTK exporter had the 6 face-center nodes partially
    permuted.  Corners and edge midpoints were correct, but ids 20--23 did not
    match VTK's actual \texttt{vtkTriQuadraticHexahedron} parametric-node
    ordering.  The visible symptom was a folded or diamond-like surface pattern
    on otherwise regular HEX27 meshes.

  \item[Decision] Correct the VTK permutation in
    \texttt{VTKCellTraits.hh}.  For the face centers, VTK expects the order
    $(x^-, x^+, y^-, y^+, z^-, z^+)$, which maps to the internal
    \texttt{fall\_n} indices $(12,14,10,16,4,22)$.

  \item[Verification] Add a regression test that compares the exporter
    permutation against the coordinates returned by
    \texttt{vtkTriQuadraticHexahedron::GetParametricCoords()}, instead of
    relying on a hand-written table alone.

  \item[Source] \texttt{src/post-processing/VTK/VTKCellTraits.hh},
    \texttt{tests/test\_simplex.cpp}.
\end{description}


% ─────────────────────────────────────────────────────────────────────────
\section{Nonlinear Showcase (\texttt{main.cpp})}
\label{sec:showcase}

The main executable exercises 7~constitutive/kinematic formulations on both
HEX27 and TET10 meshes (14~cases total) on a cantilever beam:

\begin{center}
\renewcommand{\arraystretch}{1.15}
\begin{tabular}{clllc}
  \toprule
  Case & Kinematics & Constitutive & Load & Steps \\
  \midrule
  1 & SmallStrain      & Linear elastic  & 0.05 & 1 (linear) \\
  2 & TotalLagrangian  & SVK             & 0.05 & 1 \\
  3 & TotalLagrangian  & SVK             & 0.50 & 5 \\
  4 & TotalLagrangian  & Neo-Hookean     & 0.50 & 5 \\
  5 & UpdatedLagrangian & SVK            & 0.05 & 1 \\
  6 & UpdatedLagrangian & Neo-Hookean    & 0.50 & 5 \\
  7 & SmallStrain      & J2 plasticity   & 1.00 & 20 \\
  \bottomrule
\end{tabular}
\end{center}

\paragraph{Material parameters:}
$E = 200$, $\nu = 0.3$, $\sigma_y = 0.250$, $H = 10.0$.

\paragraph{Geometry:}
$[0, 10] \times [0, 0.4] \times [0, 0.8]$.
Fixed face at $x = 0$, traction on face at $x = 10$.

\paragraph{Meshes:}
\begin{itemize}[nosep]
  \item \texttt{data/input/box.msh}: HEX27 (20~elements, 315~nodes).
  \item \texttt{data/input/box\_tet10.msh}: TET10 ($\sim$1329~elements,
        2840~nodes).
\end{itemize}

\paragraph{Post-processing:}
All cases use \texttt{compute\_material\_fields()} with adaptive dispatch
(lumped $L^2$ when safe, polynomial patch recovery otherwise), producing dual
VTU outputs (mesh + Gauss cloud) in
\texttt{data/output/}.


% ─────────────────────────────────────────────────────────────────────────
\section{Test Suite}
\label{sec:tests}

The simplex-related tests in \texttt{tests/test\_simplex.cpp} include 49~test
cases covering:
\begin{itemize}[nosep]
  \item Node count, reference coordinates, and shape function values for
        TET4 and TET10.
  \item Barycentric coordinate partition of unity and non-negativity.
  \item Shape function gradient consistency (finite difference vs.\ analytical).
  \item Face topology and local-to-global node mapping.
  \item Quadrature weight sums ($= 1/D!$ for the reference simplex).
  \item All-positive-weight verification for HMS and conical product rules.
  \item Points-inside-simplex verification for conical product.
  \item Polynomial exactness tests (degrees 3, 5, 7) for conical product
        integrators.
  \item \texttt{ConicalProductIntegrator<D, N>} template instantiation for
        $D \in \{1, 2, 3\}$, $N \in \{1, 2, 3, 4\}$.
\end{itemize}

All 49~tests and 16~element geometry tests pass as of the current revision.


% ─────────────────────────────────────────────────────────────────────────
\section{File Inventory}
\label{sec:file_inventory}

\begin{center}
\small
\renewcommand{\arraystretch}{1.1}
\begin{tabular}{p{6.8cm}p{7.5cm}}
  \toprule
  File & Purpose \\
  \midrule
  \texttt{src/geometry/SimplexCell.hh}
    & Reference simplex: coords, basis, topology (consteval) \\
  \texttt{src/elements/element\_geometry/SimplexElement.hh}
    & Element type combining geometry + integrator \\
  \texttt{src/numerics/.../SimplexQuadrature.hh}
    & Classical symmetric rules (HMS, Keast) \\
  \texttt{src/numerics/.../StroudConicalProduct.hh}
    & Golub--Welsch + Duffy + ConicalProductIntegrator \\
  \texttt{src/mesh/gmsh/GmshDomainBuilder.hh}
    & Gmsh import with TET10 node reordering \\
  \texttt{src/post-processing/VTK/VTKModelExporter.hh}
    & L2 and SPR projection, VTU export \\
  \texttt{data/input/box\_tet.geo}
    & Gmsh geometry for tetrahedral cantilever mesh \\
  \texttt{data/input/box\_tet10.msh}
    & TET10 mesh (2840 nodes) \\
  \texttt{tests/test\_simplex.cpp}
    & 49 simplex + conical product unit tests \\
  \texttt{main.cpp}
    & 14-case nonlinear showcase \\
  \bottomrule
\end{tabular}
\end{center}

% =============================================================================
%  Chapter 7 — Element Architecture: Design Patterns and Extensibility
% =============================================================================

\chapter{Element Architecture: Design Patterns and Extensibility}
\label{ch:architecture}

This chapter documents the layered element architecture of \texttt{fall\_n},
focusing on two central abstractions --- \texttt{ElementGeometry} and
\texttt{FEM\_Element} --- their interaction with policy-based design, and the
resulting customisation points.  We compare the approach against the
architectures of OpenSees~\cite{OpenSees} and Kratos
Multiphysics~\cite{Kratos}, highlighting trade-offs in extensibility, compile
time, and runtime efficiency.


% ═══════════════════════════════════════════════════════════════════════════
\section{Guiding Principles}
\label{sec:principles}

\texttt{fall\_n} is designed around three principles that permeate the element
subsystem:

\begin{enumerate}[nosep]
  \item \textbf{Zero-cost abstractions.}  Template instantiation resolves
        element type, quadrature rule, kinematic formulation, and constitutive
        model at compile time.  No virtual dispatch in the inner quadrature
        loop.  The compiler sees the full call chain from Gauss-point
        integration down to the shape function evaluation, enabling inlining
        and SIMD vectorisation.

  \item \textbf{Separation of geometry from physics.}  The geometric
        interpolation (shape functions, Jacobian, quadrature) is decoupled
        from the physical element (strain--displacement operator, constitutive
        law, DOF assembly).  The same \texttt{ElementGeometry<3>} serves both
        a \texttt{ContinuumElement} (3D solid) and a \texttt{SurfaceLoad}
        (boundary traction).

  \item \textbf{Type erasure at system boundaries only.}  Within a single
        element type, everything is monomorphic and stack-allocated.  Virtual
        dispatch is introduced only where heterogeneity is required: storing
        mixed element types in a single container (\texttt{FEM\_Element}) or
        interfacing geometry with the mesh manager
        (\texttt{ElementGeometry::pimpl\_}).
\end{enumerate}

These principles are realised through two complementary patterns:
\emph{policy-based design} (compile-time customisation) and
\emph{external polymorphism} (runtime type erasure).


% ═══════════════════════════════════════════════════════════════════════════
\section{Layer Decomposition}
\label{sec:layers}

The element subsystem is organised in four layers, each with a clear
responsibility boundary:

\begin{center}
\begin{tikzpicture}[
    layer/.style={draw, rounded corners=3pt, minimum width=13cm,
                  minimum height=1.1cm, font=\small},
    arr/.style={-{Stealth[length=5pt]}, thick}
  ]
  \node[layer, fill=blue!8]   (L4) at (0, 4.5)
    {\textbf{Layer 4 — Analysis} : \texttt{Model}, \texttt{LinearAnalysis},
     \texttt{NLAnalysis}, \texttt{DynamicAnalysis}};
  \node[layer, fill=green!8]  (L3) at (0, 3.0)
    {\textbf{Layer 3 — Physical Element} : \texttt{ContinuumElement},
     \texttt{BeamElement}, \texttt{SurfaceLoad}};
  \node[layer, fill=orange!8] (L2) at (0, 1.5)
    {\textbf{Layer 2 — Element Geometry} : \texttt{ElementGeometry<dim>}
     (type-erased handle)};
  \node[layer, fill=red!8]    (L1) at (0, 0.0)
    {\textbf{Layer 1 — Concrete Geometry} : \texttt{LagrangeElement<Dim,N...>},
     \texttt{SimplexElement<Dim,Top,Ord>}};

  \draw[arr] (L1) -- (L2);
  \draw[arr] (L2) -- (L3);
  \draw[arr] (L3) -- (L4);
\end{tikzpicture}
\end{center}

\subsection{Layer~1: Concrete Geometry Types}

These are pure value types with \emph{no virtual functions}.  They encode
the reference element, shape functions, and Jacobian evaluation as
compile-time-resolved operations.

\paragraph{LagrangeElement\texttt{<Dim, N...>}.}
A tensor-product Lagrange element parameterised by the embedding dimension
\texttt{Dim} and per-axis node counts \texttt{N...} (variadic).

\begin{center}
\small
\begin{tabular}{lcl}
  \toprule
  Instantiation & Topology & Gmsh type \\
  \midrule
  \texttt{LagrangeElement<3, 2, 2, 2>} & HEX8  & 5 \\
  \texttt{LagrangeElement<3, 3, 3, 3>} & HEX27 & 12 \\
  \texttt{LagrangeElement<3, 2, 2>}    & QUAD4 (surface) & 3 \\
  \texttt{LagrangeElement<3, 3, 3>}    & QUAD9 (surface) & 10 \\
  \texttt{LagrangeElement<3, 2>}       & LINE2 (edge)    & 1 \\
  \bottomrule
\end{tabular}
\end{center}

Key design features:
\begin{itemize}[nosep]
  \item The reference cell (\texttt{geometry::cell::LagrangianCell<N...>}) is a
        \texttt{constexpr} object holding the nodal coordinates, the
        tensor-product basis, and the subentity topology --- all computed at
        compile time.
  \item The topological dimension is deduced from the parameter pack:
        $\texttt{sizeof...(N)}$.  Thus a surface element is simply a
        \texttt{LagrangeElement<3, 3, 3>} --- same code path, fewer
        template parameters.
  \item Shape functions are products of 1D Lagrange polynomials, evaluated by
        index decomposition.  No lookup table; the compiler folds constant
        expressions.
  \item The \texttt{local\_index\_[]} array enables node reordering (e.g., Gmsh
        to internal convention) without modifying the connectivity array.
\end{itemize}

\paragraph{SimplexElement\texttt{<Dim, TopDim, Order>}.}
A simplex element parameterised by embedding dimension, topological dimension,
and polynomial order.  It mirrors the \texttt{LagrangeElement} public
interface but derives its reference cell from
\texttt{geometry::simplex::SimplexCell<TopDim, Order>}, which stores the
barycentric-based Lagrange shape functions for simplicial geometry.

Both families share a structurally identical interface (node access,
\texttt{H()}, \texttt{dH\_dx()}, \texttt{evaluate\_jacobian()},
\texttt{map\_local\_point()}) but they are \emph{unrelated types} ---
there is no inheritance hierarchy linking them.  Uniformity is achieved
through duck-typing in the next layer.


\subsection{Layer~2: ElementGeometry (Type-Erased Handle)}
\label{sec:elem_geom}

\texttt{ElementGeometry<dim>} is a type-erased value-semantic wrapper that
stores \emph{any} concrete geometry type together with its quadrature rule
behind a single, uniform interface.  It is the pivot of the entire element
subsystem.

\paragraph{Pattern: External Polymorphism.}
The design follows the \emph{External Polymorphism} pattern
(Cleary et al.\ 1997, popularised by Sean Parent's ``Inheritance Is the Base
Class of Evil'' talk):

\begin{enumerate}[nosep]
  \item \textbf{Concept} (\texttt{impl::ElementGeometryConcept<dim>}) ---
        abstract base with pure virtual methods: \texttt{H()},
        \texttt{dH\_dx()}, \texttt{evaluate\_jacobian()},
        \texttt{differential\_measure()}, \texttt{num\_integration\_points()},
        quadrature accessors, and subentity topology.

  \item \textbf{Model} (\texttt{impl::OwningModel\_ElementGeometry%
        <ElementType, IntegrationStrategy>}) ---
        template class deriving from the concept, \emph{owning} a concrete
        \texttt{ElementType} and \texttt{IntegrationStrategy} by value.
        Each virtual override delegates to the stored objects.

  \item \textbf{Handle} (\texttt{ElementGeometry<dim>}) ---
        holds a \texttt{std::unique\_ptr<ElementGeometryConcept<dim>>}
        (\texttt{pimpl\_}).  Provides value semantics via \texttt{clone()},
        move defaulted.
\end{enumerate}

\begin{center}
\begin{tikzpicture}[
    box/.style={draw, rounded corners=2pt, minimum width=4cm,
                minimum height=0.9cm, font=\small\ttfamily},
    arr/.style={-{Stealth[length=4pt]}, thick},
    impl/.style={-{Stealth[length=4pt]}, thick, dashed}
  ]
  \node[box, fill=orange!10] (handle) at (0, 3)
    {ElementGeometry<dim>};
  \node[box, fill=gray!10]   (concept) at (0, 1.5)
    {ElementGeometryConcept<dim>};
  \node[box, fill=green!10]  (model1) at (-3.5, 0)
    {OwningModel<Lag,GL>};
  \node[box, fill=green!10]  (model2) at (3.5, 0)
    {OwningModel<Simp,CP>};

  \draw[impl] (handle) -- node[right, font=\scriptsize]
    {unique\_ptr} (concept);
  \draw[arr]  (model1) -- (concept);
  \draw[arr]  (model2) -- (concept);

  \node[font=\scriptsize, below=2pt of model1]
    {LagrangeElement + GaussLegendre};
  \node[font=\scriptsize, below=2pt of model2]
    {SimplexElement + ConicalProduct};
\end{tikzpicture}
\end{center}

\paragraph{Why External Polymorphism and not classical inheritance?}

\begin{itemize}[nosep]
  \item \texttt{LagrangeElement} and \texttt{SimplexElement} remain pure value
        types --- no vtable pointer, no heap allocation, fully
        \texttt{constexpr}-eligible.  They can be used independently of the
        type-erasure layer (e.g., in constexpr unit tests).
  \item The virtual dispatch lives in the \emph{wrapper}, not in the element
        itself.  Concrete element code retains full optimisation potential.
  \item New element families (e.g., NURBS, serendipity, spectral) are added by
        writing a new value type with the required interface --- no base class
        to inherit from, no code changes in existing elements.
\end{itemize}

\paragraph{Differential measure dispatch.}
The \texttt{differential\_measure()} method in \texttt{OwningModel} uses
\texttt{if constexpr} on the compile-time-known topological dimension:
\begin{lstlisting}[language=C++]
if constexpr (topological_dim == dim)
    return std::abs(J.determinant());        // volume
else if constexpr (topological_dim == 1)
    return J.col(0).norm();                  // arc length
else if constexpr (topological_dim == 2 && dim == 3)
    return J.col(0).cross(J.col(1)).norm();  // surface area
\end{lstlisting}
Despite crossing the virtual boundary, the Jacobian matrix itself is
stack-allocated (Eigen fixed-size \texttt{Matrix<double, dim, Dynamic, \ldots,
dim, dim>}), so no heap allocation occurs inside the hot loop.

\paragraph{Integration via template callback.}
\texttt{ElementGeometry} exposes a non-virtual template method
\texttt{integrate(f)} that accepts any callable \texttt{f(LocalCoordView)}
and loops over quadrature points:
\begin{lstlisting}[language=C++]
template<std::invocable<LocalCoordView> F>
auto integrate(F&& f) const {
    // ... sum w_i * |J(xi_i)| * f(xi_i) ...
}
\end{lstlisting}
Because \texttt{f} is a template parameter (not \texttt{std::function}), the
compiler inlines the integrand.  Only the quadrature point data (coordinates,
weights) and the Jacobian evaluation pass through the virtual boundary; the
user's element-level computation stays in the monomorphic world.


\subsection{Layer~3: Physical Elements}
\label{sec:physical_elements}

Physical elements combine an \texttt{ElementGeometry} with material data and
DOF assembly logic.  Each physical element type is a class template
parameterised by a \emph{MaterialPolicy} and a \emph{KinematicPolicy}.

\paragraph{ContinuumElement\texttt{<MaterialPolicy, ndof, KinematicPolicy>}.}
The workhorse for 1D/2D/3D solid mechanics:
\begin{itemize}[nosep]
  \item Holds a non-owning pointer to an \texttt{ElementGeometry<dim>} (owned
        by the \texttt{Domain}).
  \item Owns the \texttt{MaterialPoint} array (one per Gauss point), each
        holding a \texttt{Material<MaterialPolicy>} with its constitutive
        strategy.
  \item Provides \texttt{inject\_K()}, \texttt{compute\_internal\_forces()},
        \texttt{inject\_tangent\_stiffness()}, \texttt{commit\_material\_state()}
        --- the full \texttt{FiniteElement} concept interface.
  \item The strain--displacement operator $\mathbf{B}$ is delegated to the
        \texttt{KinematicPolicy} (e.g., \texttt{SmallStrain} or
        \texttt{TotalLagrangian}), resolved at compile time.
\end{itemize}

\paragraph{BeamElement\texttt{<BeamPolicy, Dim, AsmPolicy>}.}
Timoshenko beam with section integration:
\begin{itemize}[nosep]
  \item Two-node, local-frame kinematics with a rotation matrix $\mathbf{R}$.
  \item Integration along the beam axis uses
        \texttt{MaterialSection<BeamPolicy>} objects instead of
        \texttt{MaterialPoint}.
  \item The \texttt{AsmPolicy} template parameter selects between direct
        assembly (\texttt{DirectAssembly}) and static condensation
        (\texttt{CondensedAssembly}) --- a compile-time strategy switch.
\end{itemize}

\paragraph{SurfaceLoad.}
Not a ``physical element'' in the structural sense, but uses the
\texttt{ElementGeometry} quadrature machinery to compute consistent nodal
forces from a distributed traction on a boundary face.


\subsection{Layer~4: Analysis Level}
\label{sec:analysis_layer}

\texttt{Model<MaterialPolicy, KinematicPolicy, ndofs, ElemPolicy>} owns the
element container, the \texttt{Domain}, and the PETSc objects
(\texttt{PetscSection}, matrices, vectors).

The \texttt{ElemPolicy} template parameter controls element storage
(\cref{sec:elem_policy}).  For homogeneous meshes, the container is a raw
\texttt{std::vector<ContinuumElement<\ldots>>} with zero overhead.  For
mixed meshes, it is \texttt{std::vector<FEM\_Element>} with one virtual
dispatch per element call.


% ═══════════════════════════════════════════════════════════════════════════
\section{FEM\_Element: Top-Level Type Erasure}
\label{sec:fem_element}

\texttt{FEM\_Element} applies the same External Polymorphism pattern as
\texttt{ElementGeometry}, but at the assembly level:

\begin{center}
\begin{tikzpicture}[
    box/.style={draw, rounded corners=2pt, minimum width=4.2cm,
                minimum height=0.9cm, font=\small\ttfamily},
    arr/.style={-{Stealth[length=4pt]}, thick},
    impl/.style={-{Stealth[length=4pt]}, thick, dashed}
  ]
  \node[box, fill=orange!10] (handle) at (0, 3)
    {FEM\_Element};
  \node[box, fill=gray!10]   (concept) at (0, 1.5)
    {Concept (inner struct)};
  \node[box, fill=green!10]  (model1) at (-3.5, 0)
    {Model<ContinuumElem>};
  \node[box, fill=green!10]  (model2) at (3.5, 0)
    {Model<BeamElement>};

  \draw[impl] (handle) -- node[right, font=\scriptsize]
    {unique\_ptr} (concept);
  \draw[arr]  (model1) -- (concept);
  \draw[arr]  (model2) -- (concept);
\end{tikzpicture}
\end{center}

Key design decisions:

\begin{itemize}[nosep]
  \item \textbf{Value semantics.} Copy via \texttt{clone()}, move defaulted.
        No shared ownership, no reference counting.
  \item \textbf{Self-wrap guard.} A \texttt{requires} clause prevents wrapping
        a \texttt{FEM\_Element} inside another \texttt{FEM\_Element} (double
        indirection).
  \item \textbf{Recursive concept satisfaction.} \texttt{FEM\_Element} itself
        satisfies the \texttt{FiniteElement} concept, verified by a
        \texttt{static\_assert} at the bottom of the header.  This means it
        can be nested inside other type-erased wrappers without losing the
        concept constraint.
  \item \textbf{Optional interface methods.} Mass matrix operations
        (\texttt{density()}, \texttt{inject\_mass()}) use \texttt{if constexpr}
        inside \texttt{Model<T>} to detect whether the concrete element
        supports them, defaulting to a no-op otherwise.  This avoids bloating
        the concept with operations not all elements provide.
\end{itemize}

There is also a \texttt{StructuralElement} wrapper, identical in pattern,
that provides type safety for structural-only models (beams, shells).  Both
derive from \emph{no common base class} --- they are independent type-erasure
handles, unified only by the \texttt{FiniteElement} concept.


% ═══════════════════════════════════════════════════════════════════════════
\section{Policy-Based Design: Customisation Points}
\label{sec:policies}

\texttt{fall\_n} uses policies --- template parameters that select
compile-time strategies --- at multiple levels.  Each policy is a
\emph{customisation point}: a seam where the user can inject different
behaviour without modifying existing code.

\subsection{MaterialPolicy}
\label{sec:mat_policy}

A \emph{trait bag} that binds kinematic and conjugate types:

\begin{center}
\small
\begin{tabular}{lllc}
  \toprule
  Policy & \texttt{StrainT} & \texttt{StressT} & \texttt{dim} \\
  \midrule
  \texttt{ThreeDimensionalMaterial} & \texttt{Strain<6>} & \texttt{Stress<6>} & 3 \\
  \texttt{PlaneMaterial}            & \texttt{Strain<3>} & \texttt{Stress<3>} & 2 \\
  \texttt{UniaxialMaterial}         & \texttt{Strain<1>} & \texttt{Stress<1>} & 1 \\
  \texttt{TimoshenkoBeam3D}         & \texttt{BeamStrain<6,3>} & \texttt{BeamForces<6>} & 3 \\
  \texttt{TimoshenkoBeam2D}         & \texttt{BeamStrain<3,2>} & \texttt{BeamForces<3>} & 2 \\
  \bottomrule
\end{tabular}
\end{center}

All generic code (\texttt{Material<P>}, \texttt{MaterialPoint<P>},
\texttt{MaterialState<P>}, \texttt{ElasticRelation<P>}) is parameterised
on a single \texttt{MaterialPolicy} and extracts its types via
\texttt{P::StrainT}, \texttt{P::StressT}, etc.

\paragraph{Extensibility.}
Adding a shell element requires only:
\begin{enumerate}[nosep]
  \item Define \texttt{ShellMaterial<N, 3>} with appropriate
        \texttt{ShellGeneralizedStrain<N,3>} and
        \texttt{ShellResultants<N>}.
  \item Write a constitutive model conforming to
        \texttt{ConstitutiveRelation<ShellMaterial<N,3>>}.
  \item Write the element class (\texttt{ShellElement}) satisfying
        \texttt{FiniteElement}.
\end{enumerate}
No existing code changes; the policy mechanism propagates types automatically.


\subsection{KinematicPolicy}
\label{sec:kin_policy}

Selects the strain--displacement relationship at compile time:

\begin{center}
\small
\begin{tabular}{lp{5cm}cc}
  \toprule
  Policy & Description & Linear? & $K_\sigma$? \\
  \midrule
  \texttt{SmallStrain}       & $\bvarepsilon = \nabla^s \mathbf{u}$ & yes & no \\
  \texttt{TotalLagrangian}   & $\mathbf{E} = \frac{1}{2}(\bF^T\bF - \bI)$ & no & yes \\
  \texttt{UpdatedLagrangian} & Cauchy strain in current config. & no & yes \\
  \texttt{Corotational}      & Large rotation, small strain & no & yes \\
  \bottomrule
\end{tabular}
\end{center}

Each policy is a struct with static methods:
\begin{lstlisting}[language=C++]
struct SmallStrain {
    static constexpr bool is_geometrically_linear    = true;
    static constexpr bool needs_geometric_stiffness  = false;
    
    template <std::size_t dim>
    static GPKinematics<dim> evaluate(
        ElementGeometry<dim>* geo, std::size_t num_nodes,
        std::size_t ndof, const Array& Xi,
        const Eigen::VectorXd& u_e);

    template <std::size_t dim>
    static auto compute_B(ElementGeometry<dim>* geo, ...);
};
\end{lstlisting}

\texttt{ContinuumElement} calls \texttt{KinematicPolicy::evaluate<dim>()} at
each quadrature point.  The boolean traits steer \texttt{if constexpr}
branches inside the element (e.g., add $K_\sigma$ only when
\texttt{needs\_geometric\_stiffness == true}).

\paragraph{Extensibility.}
A new formulation (e.g., logarithmic strain) is added by defining a new
struct satisfying \texttt{KinematicPolicyConcept} with the required static
methods and boolean traits.  No changes to \texttt{ContinuumElement} or any
analysis class.


\subsection{ElementPolicy (Storage Strategy)}
\label{sec:elem_policy}

Controls how \texttt{Model} stores its element collection:

\begin{center}
\small
\begin{tabular}{lp{4.5cm}cc}
  \toprule
  Policy & Container & Homogeneous & Overhead \\
  \midrule
  \texttt{SingleElementPolicy<E>} & \texttt{std::vector<E>}
    & yes & none \\
  \texttt{MultiElementPolicy}     & \texttt{std::vector<FEM\_Element>}
    & no  & 1 vtable + 1 ptr chase \\
  \bottomrule
\end{tabular}
\end{center}

The concept \texttt{ElementPolicyLike} enforces the policy interface:
\begin{lstlisting}[language=C++]
template <typename P>
concept ElementPolicyLike = requires {
    typename P::element_type;
    typename P::container_type;
    { P::is_homogeneous } -> std::convertible_to<bool>;
} && FiniteElement<typename P::element_type>;
\end{lstlisting}

\paragraph{Default path.}
$\texttt{Model<ThreeDimensionalMaterial>}$ defaults to:\\
$\texttt{SingleElementPolicy<ContinuumElement<ThreeDimensionalMaterial, 3,}$
$\texttt{SmallStrain>>}$,\\
yielding a contiguous \texttt{std::vector} of monomorphic elements --- optimal
for cache locality and auto-vectorisation.


\subsection{AssemblyPolicy (Beams)}

Beam elements accept an \texttt{AsmPolicy} parameter:
\begin{itemize}[nosep]
  \item \texttt{DirectAssembly} --- full stiffness matrix, no condensation.
  \item \texttt{CondensedAssembly} --- static condensation of internal DOFs
        before injection into PETSc.
\end{itemize}
Both are stateless tag types; the choice is resolved at compile time via
\texttt{if constexpr} inside \texttt{BeamElement::inject\_K()}.

\subsection{ConstitutiveStrategy (Material Update)}

The material update mechanism (\texttt{Material<P>}) delegates to a
type-erased strategy object (\texttt{IntegrationStrategy<P>}), enabling
runtime polymorphism over the constitutive response:
\begin{itemize}[nosep]
  \item \texttt{ElasticRelation<P>} --- linear elastic.
  \item \texttt{HyperelasticRelation<P>} --- Neo-Hookean, SVK, etc.
  \item \texttt{J2Plasticity<P>} --- return mapping with isotropic hardening.
\end{itemize}
This is the \emph{only} place in the element quadrature loop where runtime
polymorphism is used, because the constitutive model is genuinely a runtime
choice (same mesh, different materials per physical group).


% ═══════════════════════════════════════════════════════════════════════════
\section{Compile-Time Cost Analysis}
\label{sec:compile_time}

Policy-based design with heavy template usage has a well-known compile-time
cost.  \texttt{fall\_n} mitigates this through several strategies:

\begin{enumerate}[nosep]
  \item \textbf{Header-only by necessity, not by choice.}  The library is
        header-only because C++ templates must be visible at the point of
        instantiation.  However, explicit template instantiation could be
        added in \texttt{.cpp} files for the most common configurations to
        reduce redundant codegen across translation units.

  \item \textbf{Type erasure as a firewall.}  \texttt{ElementGeometry<dim>}
        and \texttt{FEM\_Element} act as \emph{template firewalls}: code that
        consumes an \texttt{ElementGeometry<3>} does not need to see the
        headers of \texttt{LagrangeElement} or \texttt{SimplexElement}.
        This significantly reduces transitive include graphs.

  \item \textbf{Consteval reference cells.}  The \texttt{LagrangianCell<N...>}
        and \texttt{SimplexCell<TopDim, Order>} objects with their basis
        functions and topological tables are \texttt{constexpr}/\texttt{consteval},
        evaluated by the compiler once.  They produce no runtime code unless
        actually called; the compiler may fold constant subexpressions entirely.

  \item \textbf{Concept constraints as documentation.}  The
        \texttt{FiniteElement}, \texttt{KinematicPolicyConcept},
        \texttt{ElementPolicyLike} concepts serve as lightweight checks that
        produce clear error messages at instantiation time, avoiding the
        traditional deep-template-error-message problem.
\end{enumerate}

\paragraph{Measured impact.}
A full rebuild of \texttt{fall\_n} (main + 12 test targets) takes approximately
105~seconds on a single core with GCC~11 and debug optimisation (\texttt{-g}).
The dominant cost is Eigen's template instantiation inside
\texttt{ContinuumElement} and \texttt{BeamElement}; the element geometry layer
adds minimal overhead because the type-erasure boundary limits template
propagation.


% ═══════════════════════════════════════════════════════════════════════════
\section{Comparison with OpenSees and Kratos}
\label{sec:comparison}

\subsection{Overview}

\begin{center}
\small
\renewcommand{\arraystretch}{1.2}
\begin{tabular}{p{3.8cm}p{3.8cm}p{3.8cm}p{3.8cm}}
  \toprule
  Aspect & \textbf{OpenSees} & \textbf{Kratos} & \textbf{fall\_n} \\
  \midrule
  Language &
    C++ (C++11/14) &
    C++ (C++17) + Python &
    C++ (C++20/23) \\
  Polymorphism model &
    Classical inheritance &
    Classical inheritance + Python bridge &
    Policies + External Polymorphism \\
  Element base &
    \texttt{Element} (virtual) &
    \texttt{Element} (virtual) &
    \texttt{FiniteElement} (concept) \\
  Geometry base &
    \texttt{CrdTransf} (virtual) &
    \texttt{Geometry<TPoint>} (virtual) &
    \texttt{ElementGeometry<dim>} (type-erased) \\
  Material base &
    \texttt{UniaxialMaterial} / \texttt{NDMaterial} (virtual) &
    \texttt{ConstitutiveLaw} (virtual) &
    \texttt{Material<Policy>} (policy + type-erased strategy) \\
  Dispatch cost &
    1--3 vtable calls per GP &
    1--2 vtable calls per GP &
    0 (homogeneous) or 1 (heterogeneous) \\
  \bottomrule
\end{tabular}
\end{center}

\subsection{OpenSees}

OpenSees~\cite{OpenSees} uses a deep inheritance hierarchy:
$\texttt{DomainComponent} \to \texttt{Element} \to$ concrete types.
Materials form an even deeper tree:
$\texttt{Material} \to \texttt{UniaxialMaterial} \to$ specific laws.

\paragraph{Strengths.}
\begin{itemize}[nosep]
  \item Extremely mature (20+ years), battle-tested in earthquake engineering.
  \item Simple runtime scripting via Tcl (later Python) --- users define
        models without recompiling.
  \item Large element library (200+ elements).
\end{itemize}

\paragraph{Weaknesses (from an architecture perspective).}
\begin{itemize}[nosep]
  \item \textbf{Monolithic virtual hierarchy.}  All elements inherit from
        \texttt{Element}, which carries 30+ virtual methods including
        \texttt{getResistingForce()}, \texttt{getTangentStiff()}, etc.
        Adding a new element requires implementing all of them (or relying
        on default no-ops that may silently produce incorrect results).
  \item \textbf{No separation of geometry and physics.}  A \texttt{Brick20N}
        element hard-codes its shape functions, Jacobian, AND constitutive
        integration in a single class.  Changing the quadrature rule requires
        modifying the element source.
  \item \textbf{Dynamic dispatch in the inner loop.}  Every call to
        \texttt{material->getStress()} and \texttt{material->getTangent()}
        at each Gauss point goes through virtual dispatch.  While the cost
        is small relative to the linear algebra, it inhibits inlining and
        auto-vectorisation across the material boundary.
  \item \textbf{No concept or constraint checking.}  Adding a new element
        with a typo in a virtual override signature compiles silently and
        produces incorrect results at runtime.
\end{itemize}

\paragraph{Contrast with fall\_n.}
\texttt{fall\_n} eliminates the monolithic base class by defining
\texttt{FiniteElement} as a C++20 concept that enforces the interface at
compile time.  The element is free to organise its internals however it
pleases; the concept checks only what the assembler needs.  Geometry is
factored out into \texttt{ElementGeometry}, so the same mesh can serve
continuum, surface-load, and post-processing operations.


\subsection{Kratos Multiphysics}

Kratos~\cite{Kratos} uses a more modular architecture than OpenSees, with
explicit separation of geometry (\texttt{Geometry<TPoint>}), element
(\texttt{Element}), condition (\texttt{Condition}), constitutive law
(\texttt{ConstitutiveLaw}), and process (\texttt{Process}).

\paragraph{Strengths.}
\begin{itemize}[nosep]
  \item \textbf{Multi-physics by design.}  The \texttt{ModelPart} /
        \texttt{ProcessInfo} / \texttt{Variable} registry enables coupling of
        disparate physics (structural, fluid, thermal) sharing the same mesh.
  \item \textbf{Python-driven workflow.}  Element selection, boundary
        conditions, solver parameters, and post-processing are all scriptable
        in Python without recompilation.
  \item \textbf{Geometry separated from element.}  \texttt{Geometry<TPoint>}
        handles shape functions and Jacobian; \texttt{Element} handles physics.
        This is conceptually analogous to \texttt{fall\_n}'s two-layer split.
\end{itemize}

\paragraph{Weaknesses (from an architecture perspective).}
\begin{itemize}[nosep]
  \item \textbf{Runtime overhead of the variable system.}
        Each nodal value is stored in a hash-map-like container
        (\texttt{Variable<T>} keys), requiring a lookup per node per variable
        access.  This sacrifices cache locality and prevents compile-time
        resolution of DOF layouts.
  \item \textbf{Virtual geometry hierarchy.}  \texttt{Geometry<TPoint>} is an
        inheritance tree: $\texttt{Geometry} \to \texttt{Triangle2D3}$, etc.
        Shape functions are declared \texttt{virtual}, adding dispatch cost at
        each Gauss point.  There is no mechanism to ``inline'' the geometry into
        the element at compile time.
  \item \textbf{Heavyweight element interface.}  \texttt{Element} carries
        methods for \texttt{CalculateLocalSystem},
        \texttt{CalculateMassMatrix}, \texttt{CalculateDampingMatrix},
        \texttt{GetDofList}, etc.  Many elements implement subsets, leaving
        the rest as no-ops or throwing exceptions --- a code-smell of
        interface bloat.
  \item \textbf{No compile-time guarantees.}  Registering a new element
        involves adding entries to a factory map at runtime.  Errors in the
        element's interface (wrong return type, missing method) are caught
        only at link time or, worse, at runtime.
\end{itemize}

\paragraph{Contrast with fall\_n.}
\texttt{fall\_n}'s policy-based design replaces the Kratos variable registry
with compile-time type parameters: the \texttt{MaterialPolicy} selects the
strain/stress types, the \texttt{KinematicPolicy} selects the B-operator, and
the \texttt{ElementPolicy} selects the storage strategy --- all resolved by
the compiler.  The trade-off is that \texttt{fall\_n} requires recompilation
to change element types, whereas Kratos allows runtime scripting.


\subsection{Summary: Architectural Trade-offs}
\label{sec:tradeoffs}

\begin{center}
\small
\renewcommand{\arraystretch}{1.2}
\begin{tabular}{p{4cm}ccc}
  \toprule
  Trade-off & OpenSees & Kratos & fall\_n \\
  \midrule
  Runtime scripting         & \checkmark\checkmark & \checkmark\checkmark\checkmark & --- \\
  Compile-time safety       & ---  & ---  & \checkmark\checkmark\checkmark \\
  Zero-cost inner loop      & ---  & ---  & \checkmark\checkmark \\
  Geometry--physics split   & ---  & \checkmark\checkmark & \checkmark\checkmark\checkmark \\
  Heterogeneous elements    & \checkmark\checkmark\checkmark & \checkmark\checkmark\checkmark & \checkmark\checkmark \\
  Multi-physics coupling    & \checkmark & \checkmark\checkmark\checkmark & --- \\
  Maturity / element count  & \checkmark\checkmark\checkmark & \checkmark\checkmark\checkmark & \checkmark \\
  \bottomrule
\end{tabular}
\end{center}

\texttt{fall\_n} occupies a different niche: it is \emph{not} a general-purpose
production FE framework, but a research library where the primary concern is
\textbf{architectural clarity and compile-time correctness} over runtime
flexibility.  The policy-based approach is intentionally opinionated: it makes
the ``easy things trivial'' (adding a new constitutive model, kinematic
formulation, or element family within the existing type system) at the cost of
making the ``hard things hard'' (mixing unrelated physics, runtime element
selection from user input).


% ═══════════════════════════════════════════════════════════════════════════
\section{Extensibility Guide: Adding a New Element Family}
\label{sec:extensibility}

This section walks through the concrete steps to add a hypothetical
``Serendipity element'' family, illustrating each customisation point.

\begin{enumerate}
  \item \textbf{Layer~1: Concrete geometry type.}
    Create \texttt{SerendipityElement<Dim, Order>}:
    \begin{itemize}[nosep]
      \item Store node coordinates, define \texttt{H()}, \texttt{dH\_dx()},
            \texttt{evaluate\_jacobian()}, \texttt{map\_local\_point()} with
            the serendipity basis functions.
      \item No base class needed.  The type must be movable and copyable.
    \end{itemize}

  \item \textbf{Layer~1: Integration strategy.}
    Create or reuse a quadrature class (\texttt{GaussLegendre2D<p>},
    \texttt{ReducedIntegration<\ldots>}, etc.) exposing
    \texttt{reference\_integration\_point(i)} and \texttt{weight(i)}.

  \item \textbf{Layer~2: Instantiate through ElementGeometry.}
    No code changes.  The constructor template:
\begin{lstlisting}[language=C++]
ElementGeometry<3> geo(
    SerendipityElement<3, 2>{tag, node_ids},
    GaussLegendre2D<3>{});
\end{lstlisting}
    automatically instantiates an
    \texttt{OwningModel<Serendipity..., GL...>} and stores it behind the
    type-erased handle.

  \item \textbf{Layer~3: Use with ContinuumElement.}
    No code changes.  \texttt{ContinuumElement} receives a pointer to the
    geometry and calls \texttt{H()}, \texttt{dH\_dx()},
    \texttt{differential\_measure()} through the virtual interface.  The
    $\mathbf{B}$ matrix and assembly are generic.

  \item \textbf{Layer~4: Mesh import.}
    Add a case in \texttt{GmshDomainBuilder.hh} for the Gmsh element type
    code, constructing \texttt{ElementGeometry} with the serendipity concrete
    type and the desired integrator.

  \item \textbf{Optional: Type-safe model.}
    If the serendipity element will be used in its own model (no mixing),
    define:
\begin{lstlisting}[language=C++]
using SerendipitySolid =
    ContinuumElement<ThreeDimensionalMaterial, 3,
                     SmallStrain>;
using SerendipityModel =
    Model<ThreeDimensionalMaterial, SmallStrain, 3,
          SingleElementPolicy<SerendipitySolid>>;
\end{lstlisting}
    This compiles to a contiguous vector with zero virtual dispatch.
\end{enumerate}

\paragraph{What does NOT need to change:}
\begin{itemize}[nosep]
  \item \texttt{ElementGeometry.hh} (new type is auto-detected by the
        templated constructor).
  \item \texttt{ContinuumElement.hh} (works through geometry pointer).
  \item \texttt{FEM\_Element.hh} (type erasure is generic).
  \item \texttt{Model.hh} (policy-parameterised).
  \item Any analysis class (\texttt{LinearAnalysis}, \texttt{NLAnalysis}, etc.).
  \item Any material or constitutive model.
\end{itemize}

This is the central extensibility promise: \textbf{new element families are
additive} --- they require only new files, not modifications to existing ones
--- conforming to the Open/Closed Principle.


% ═══════════════════════════════════════════════════════════════════════════
\section{Runtime Efficiency Considerations}
\label{sec:efficiency}

\subsection{Virtual Dispatch Budget}

In the homogeneous path (\texttt{SingleElementPolicy}), the element assembly
loop has \textbf{zero} virtual calls for the physics: all
\texttt{inject\_K()}, \texttt{B()}, \texttt{H()}, \texttt{dH\_dx()} calls are
monomorphic and inlineable.

The only virtual calls per Gauss point are:
\begin{itemize}[nosep]
  \item \texttt{ElementGeometry::evaluate\_jacobian()} --- 1 vcall for the
        Jacobian (through pimpl).
  \item \texttt{ElementGeometry::H()}, \texttt{dH\_dx()} --- $n_\text{nodes}
        \times (1 + \text{dim})$ vcalls for shape functions and derivatives.
  \item \texttt{Material::compute\_response()} --- 1 vcall for the
        constitutive model (through the strategy pointer).
\end{itemize}
This is \emph{identical per-element overhead} to OpenSees or Kratos, but
\texttt{fall\_n} avoids the additional vtable dispatch that those frameworks
incur at the element container level.

\subsection{Cache-Friendly Layout}

\texttt{SingleElementPolicy<E>} stores all elements contiguously in a single
\texttt{std::vector<E>}.  Each element's \texttt{MaterialPoint} array is
heap-allocated (as a \texttt{std::vector<MaterialPointT>}), but its entries
are contiguous for that element.

\texttt{MultiElementPolicy} stores \texttt{FEM\_Element} objects contiguously
(the \texttt{unique\_ptr<Concept>} wrappers are packed), but the pointed-to
\texttt{Model<T>} objects are heap-allocated with potentially scattered
addresses.  This adds one pointer chase per element --- acceptable for
heterogeneous cases where the element count is moderate.

\subsection{DOF Index Caching}

\texttt{ContinuumElement} caches its flat \texttt{dof\_indices\_[]} array
(one entry per element DOF, combining all nodes).  The cache is populated
lazily on first use and reused across \texttt{inject\_K()},
\texttt{compute\_internal\_forces()}, and \texttt{inject\_tangent\_stiffness()}.
This avoids the per-call overhead of traversing the node pointers to
collect DOF indices at every assembly step.
\begin{lstlisting}[language=C++]
void ensure_dof_cache() noexcept {
    if (dofs_cached_) return;
    collect_dof_indices();
}
\end{lstlisting}

\subsection{Stack-Allocated Jacobians}

The Jacobian matrix crossing the virtual boundary is declared as:
\begin{lstlisting}[language=C++]
using JacobianMatrix = Eigen::Matrix<double, dim,
    Eigen::Dynamic, Eigen::ColMajor, dim, dim>;
\end{lstlisting}
The \texttt{MaxRows = dim}, \texttt{MaxCols = dim} template parameters tell
Eigen to use a fixed-size stack buffer (up to $3 \times 3 = 72$ bytes) with
a dynamic column count.  Since \texttt{dim} $\le 3$, no heap allocation
occurs, even though the column count is ``dynamic'' (it varies between volume
and surface elements).

This ensures that the virtual call to \texttt{evaluate\_jacobian()} returns by
value without any \texttt{new}/\texttt{delete}, making the hot path
allocation-free.

% =============================================================================
%  ch08_boundary_conditions.tex — Generic Boundary \& Initial Condition Design
% =============================================================================
\chapter{Generic Boundary and Initial Condition Infrastructure}
\label{ch:boundary-conditions}

This chapter documents the design and implementation of the generic
boundary condition (BC) and initial condition (IC) infrastructure for
\texttt{fall\_n}.  The goals are:
\begin{enumerate}[nosep]
  \item \textbf{Per-DOF granularity:}  constrain individual degrees of freedom
        (roller supports, symmetry BCs) rather than all DOFs simultaneously.
  \item \textbf{Imposed (non-zero) Dirichlet:}  prescribe non-zero
        displacement, temperature, or any nodal field value.
  \item \textbf{Compile-time node selectors:}  apply constraints via
        geometric predicates (plane, box, sphere, custom lambda) that
        resolve at compile time with zero runtime overhead.
  \item \textbf{State transfer:}  capture the complete post-solve state
        (\texttt{ModelState}) and use it as the initial condition for
        a subsequent analysis---enabling staged construction, restart,
        and multi-physics chaining.
  \item \textbf{Extensibility:}  the design should accommodate future
        field types (temperature, initial stress, pore pressure) without
        modifying the selector or constraint machinery.
\end{enumerate}


% ═════════════════════════════════════════════════════════════════════════
\section{Gap Analysis: the Legacy API}
\label{sec:bc-gap}

The original Model API exposed these methods for boundary conditions:

\begin{lstlisting}
void fix_node(size_t node_idx);                          // ALL DOFs = 0
void fix_x(double val, double tol = 1e-6);               // all DOFs at x=val
void apply_node_force(size_t node_idx, auto... f);        // nodal Neumann
void apply_surface_traction(string group, auto... t);     // consistent traction
\end{lstlisting}

\begin{table}[ht]
\centering\small
\caption{Limitations of the legacy BC interface.}
\label{tab:bc-gaps}
\begin{tabular}{lp{8cm}}
\toprule
\textbf{Limitation} & \textbf{Impact} \\
\midrule
\texttt{fix\_node} constrains \emph{all} DOFs &
Cannot model roller supports, symmetry planes, or mixed BCs. \\
\texttt{fix\_x/y/z} misnamed &
These fix all three displacement components, not just the named one.
Misleading to users.\\
No non-zero Dirichlet via Model API &
\texttt{global\_imposed\_solution\_} was created and zeroed but never
populated through the Model interface.  Only
\texttt{DynamicAnalysis::PrescribedBC} filled it.\\
\texttt{ElementInitialStress} unused &
The struct existed in \texttt{BoundaryCondition.hh} and was storable in
\texttt{BoundaryConditionSet}, but no evaluator integrated it into the
assembly pipeline.\\
Procedural selection &
Every \texttt{fix\_*} method hard-coded a coordinate loop.  No reusable
geometric predicate.\\
\texttt{VecAssemblyBegin/End} per call &
\texttt{apply\_node\_force} triggered PETSc assembly for each node. \\
\bottomrule
\end{tabular}
\end{table}


% ═════════════════════════════════════════════════════════════════════════
\section{Design Decisions and Trade-offs}
\label{sec:bc-design}

\subsection{Compile-Time Polymorphism via Concepts}

Consistent with \texttt{fall\_n}'s policy-based architecture
(\Cref{ch:architecture}), the node selector system uses C++20 concepts
rather than virtual dispatch:

\begin{lstlisting}
namespace fn {
template <typename S, typename NodeT>
concept NodeSelector = requires(const S& s, const NodeT& n) {
    { s(n) } -> std::convertible_to<bool>;
};
}
\end{lstlisting}

\noindent
Any callable \lstinline{f(const Node<dim>&) -> bool} satisfies
\texttt{NodeSelector}: lambdas, function objects, and the built-in
selectors all qualify.  The compiler resolves the predicate at compile
time, and after inlining a \texttt{PlaneSelector} collapses to a single
\texttt{fabs} + comparison---zero overhead compared to a hand-written
loop.

\paragraph{Why not \texttt{std::function}?}
A \texttt{std::function<bool(const Node\&)>} would erase the type and
force heap allocation for lambdas that exceed the small-buffer size (typically
16--32 bytes).  For selectors applied to every node in the mesh, this
per-call overhead is measurable.  The concept-based approach pays nothing
at runtime.

\paragraph{Why not virtual?}
A virtual \texttt{AbstractSelector::operator()} would add one indirect
call per node per constraint pass.  With 10\,000 nodes this is
$\sim$10\,$\mu$s (branch-predictor penalty), acceptable but unnecessary
when the concrete type is known at compile time.


\subsection{Constraint Storage: Extending the Existing Map}

The existing \texttt{ConstraintDofInfo} type already maps sieve point IDs
to (DOF index list, prescribed value list):

\begin{lstlisting}
using ConstraintDofInfo =
  std::map<PetscInt,
           std::pair<std::vector<PetscInt>,
                     std::vector<PetscScalar>>>;
\end{lstlisting}

\noindent
The new API \emph{extends} this map rather than replacing it.  The key
operations are:

\begin{enumerate}[nosep]
  \item \textbf{Insert in sorted order.}  PETSc's
        \texttt{PetscSectionSetConstraintIndices} requires increasing
        DOF indices.  The \texttt{constrain\_dof} method uses
        \texttt{std::lower\_bound} to maintain this invariant.
  \item \textbf{Merge, not overwrite.}  If DOF~0 of a node is already
        constrained and the user constrains DOF~2, both are stored.
        An update to DOF~0 changes the prescribed value without
        duplicating the index.
\end{enumerate}

\subsection{Non-Zero Dirichlet: Where Values Flow}
\label{sec:bc-nonzero}

Non-zero prescribed values are stored in \texttt{constraints\_} alongside
the DOF indices.  During \texttt{setup()}, a new
\texttt{fill\_imposed\_solution()} step extracts these values from the
constraint map into the PETSc local vector
\texttt{global\_imposed\_solution\_}:

\begin{lstlisting}
void setup() {
    setup_boundary_conditions();  // PETSc section
    setup_vectors();              // create Vecs
    fill_imposed_solution();      // NEW: u_D -> imposed
    setup_matrix();               // create Mat
}
\end{lstlisting}

\begin{remark}
Non-zero Dirichlet requires that the analysis accounts for the coupling
$\mathbf{K}_{fc}\,\mathbf{u}_c$ in the right-hand side.
\texttt{NonlinearAnalysis} handles this naturally: the internal forces
$\mathbf{f}_{\text{int}}(\mathbf{u}_{\text{free}} + \mathbf{u}_D)$ include
the prescribed part.  \texttt{LinearAnalysis} currently does \emph{not}
apply this correction---for non-zero Dirichlet, users should solve with
\texttt{NonlinearAnalysis} (which converges in one Newton step for linear
problems).
\end{remark}


\subsection{ModelState: Platform-Independent Snapshot}

For state transfer between models, we need a data structure that:
\begin{itemize}[nosep]
  \item Does \emph{not} depend on PETSc (no \texttt{Vec}, \texttt{Mat}).
  \item Does \emph{not} depend on the \texttt{MaterialPolicy} template parameter
        (so an elastic model's state can seed a plastic model).
  \item Is copyable, movable, and (future) serialisable.
\end{itemize}

\texttt{ModelState<dim>} satisfies all three.  It stores:
\begin{itemize}[nosep]
  \item Per-node: ID, displacement, velocity.
  \item Per-element, per-Gauss-point: stress and strain in Voigt notation
        (\texttt{std::array<double, nvoigt>}).
\end{itemize}

The PETSc-to-array extraction in \texttt{capture\_state()} uses
\texttt{VecGetArrayRead} (zero-copy read) and the element's
\texttt{material\_points()} accessor which recomputes stress from the
committed strain via the constitutive relation.


% ═════════════════════════════════════════════════════════════════════════
\section{Implementation}
\label{sec:bc-impl}

\subsection{NodeSelector.hh — Built-in Selectors}

\begin{table}[ht]
\centering\small
\caption{Built-in node selectors.}
\label{tab:selectors}
\begin{tabular}{lll}
\toprule
\textbf{Selector} & \textbf{Predicate} & \textbf{Parameters} \\
\midrule
\texttt{PlaneSelector<dim>}  & $|x_i - v| < \varepsilon$
                             & axis $i$, value $v$, tolerance $\varepsilon$ \\
\texttt{BoxSelector<dim>}    & $x \in [\text{lo},\text{hi}]$
                             & corners, tolerance \\
\texttt{SphereSelector<dim>} & $\|x - c\| \le r$
                             & centre $c$, radius $r$, tolerance \\
\texttt{NodeIdSelector}      & $\text{id} \in S$
                             & set of node IDs \\
\bottomrule
\end{tabular}
\end{table}

\noindent
Combinators turn simple selectors into compound predicates:

\begin{lstlisting}
auto corner = fn::select_and(
    fn::PlaneSelector<3>{0, 0.0},   // x = 0
    fn::PlaneSelector<3>{1, 0.0}    // y = 0
);
auto not_x0 = fn::select_not(fn::PlaneSelector<3>{0, 0.0});
auto union_ = fn::select_or(selectorA, selectorB);
\end{lstlisting}

\noindent
Convenience factories (\texttt{fn::on\_plane}, \texttt{fn::in\_box},
\texttt{fn::in\_sphere}, \texttt{fn::node\_ids}) provide a terser syntax.


\subsection{Model.hh — New Constraint Methods}

\begin{table}[ht]
\centering\small
\caption{New Model constraint API.}
\label{tab:constraint-api}
\begin{tabular}{lp{7.5cm}}
\toprule
\textbf{Method} & \textbf{Description} \\
\midrule
\texttt{constrain\_dof(n, d, v)}
  & Constrain DOF~$d$ of node~$n$ to value~$v$.  Sorted insertion. \\
\texttt{constrain\_node(n, vals)}
  & Constrain all DOFs of node~$n$ to the given array of values. \\
\texttt{constrain\_dof\_where(sel, d, v)}
  & Constrain DOF~$d$ at every node matching selector \texttt{sel}. \\
\texttt{constrain\_where(sel)}
  & Constrain all DOFs (homogeneous) at every matching node. \\
\texttt{constrain\_where(sel, vals)}
  & Constrain all DOFs with prescribed values at every matching node. \\
\midrule
\texttt{is\_constrained(n)}
  & Query: is any DOF of node~$n$ constrained? \\
\texttt{is\_dof\_constrained(n, d)}
  & Query: is DOF~$d$ of node~$n$ constrained? \\
\texttt{prescribed\_value(n, d)}
  & Get the prescribed displacement at DOF~$d$ (0 if free). \\
\texttt{num\_constraints()}
  & Total number of constrained DOFs across all nodes. \\
\midrule
\texttt{capture\_state()}
  & Extract current solution into a \texttt{ModelState<dim>}. \\
\texttt{apply\_initial\_displacement(s)}
  & Populate imposed solution from a \texttt{ModelState<dim>}. \\
\bottomrule
\end{tabular}
\end{table}


\subsection{Backward Compatibility}

The legacy methods \texttt{fix\_node}, \texttt{fix\_x}, \texttt{fix\_y},
\texttt{fix\_z}, and \texttt{fix\_orthogonal\_plane} remain unchanged in
semantics---they continue to constrain \emph{all} DOFs of every matching
node.  Internally, \texttt{fix\_orthogonal\_plane} now delegates to:
\begin{lstlisting}
void fix_orthogonal_plane(int i, double val, double tol) {
    constrain_where(fn::PlaneSelector<dim>{i, val, tol});
}
\end{lstlisting}
The code path is identical; only the implementation reuses the selector
infrastructure.


% ═════════════════════════════════════════════════════════════════════════
\section{Use-Case Walkthrough: Roller vs.\ Clamped Cantilever}
\label{sec:bc-roller}

Consider a cantilever beam $[0,L]\times[0,b]\times[0,h]$ loaded at
$x = L$ by a surface traction $\mathbf{t} = (0, t_y, t_z)$.

\paragraph{Clamped support (legacy):}
\begin{lstlisting}
M.fix_x(0.0);   // constrains u_x, u_y, u_z at x=0
\end{lstlisting}
Equivalent new API:
\begin{lstlisting}
M.constrain_where(fn::PlaneSelector<3>{0, 0.0});
\end{lstlisting}

\paragraph{Roller support (new API):}
Only $u_x = 0$ at $x = 0$.  Rigid body prevention requires additionally
constraining $u_y$ and $u_z$ at isolated points:
\begin{lstlisting}
// u_x = 0 on entire face
M.constrain_dof_where(fn::PlaneSelector<3>{0, 0.0}, 0, 0.0);

// Prevent rigid translation + rotation about x:
auto origin = fn::select_and(fn::PlaneSelector<3>{0, 0.0},
    fn::select_and(fn::PlaneSelector<3>{1, 0.0},
                   fn::PlaneSelector<3>{2, 0.0}));
M.constrain_dof_where(origin, 1, 0.0);  // u_y at origin
M.constrain_dof_where(origin, 2, 0.0);  // u_z at origin

auto top_z = fn::select_and(fn::PlaneSelector<3>{0, 0.0},
    fn::select_and(fn::PlaneSelector<3>{1, 0.0},
                   fn::PlaneSelector<3>{2, h}));
M.constrain_dof_where(top_z, 1, 0.0);   // prevents Rx
\end{lstlisting}

The roller case has 30~constrained DOFs (28~nodes $\times$ 1 + 2~extra)
versus 84 (28 $\times$ 3) for the clamped case.  With fewer constraints,
the beam is more flexible:

\begin{table}[ht]
\centering\small
\caption{Roller vs.\ clamped comparison (HEX8 beam, $E=200$, $\nu=0.3$).}
\label{tab:roller-clamped}
\begin{tabular}{lccc}
\toprule
\textbf{BC type} & \textbf{Constrained DOFs} & $\max|u|$ & \textbf{Ratio} \\
\midrule
Clamped & 84 & 0.594 & 1.00 \\
Roller ($u_x$ only) & 30 & 2.762 & 4.64$\times$ \\
\bottomrule
\end{tabular}
\end{table}

\noindent
This confirms: the per-DOF mechanism works correctly and produces a
physically meaningful stiffer/more-flexible response.


% ═════════════════════════════════════════════════════════════════════════
\section{State Capture and Continuation Analysis}
\label{sec:bc-state}

\subsection{ModelState Structure}

\begin{lstlisting}
template <size_t dim>
struct ModelState {
    static constexpr size_t nvoigt = dim * (dim + 1) / 2;

    struct NodalData {
        size_t node_id;
        array<double, dim> displacement;
        array<double, dim> velocity;       // zero for static
    };

    struct GaussPointData {
        array<double, nvoigt> stress;
        array<double, nvoigt> strain;
    };

    struct ElementData {
        size_t element_index;
        vector<GaussPointData> gauss_points;
    };

    vector<NodalData>   nodes;
    vector<ElementData> elements;
};
\end{lstlisting}

\subsection{Capture Workflow}

\texttt{Model::capture\_state()} extracts:
\begin{enumerate}[nosep]
  \item \textbf{Nodal displacements:}  from the PETSc local vector
        \texttt{current\_state\_} via \texttt{VecGetArrayRead} (zero-copy).
        Each node's DOF offsets from \texttt{node.dof\_index()} are used to
        index into the raw array.
  \item \textbf{Element stress/strain:}  from each \texttt{MaterialPoint}'s
        committed strain (\texttt{current\_state()}) and recomputed stress
        (\texttt{compute\_response()}).  The Voigt vectors are copied
        component-by-component into \texttt{std::array<double, nvoigt>},
        guarded by the compile-time \texttt{num\_components} constant.
\end{enumerate}

\subsection{Future: Initial Stress Integration}

The existing \texttt{ElementInitialStress<dim>} struct in
\texttt{BoundaryCondition.hh} can be populated from
\texttt{ModelState::ElementData}:

\begin{lstlisting}
for (const auto& ed : state.elements) {
    ElementInitialStress<dim> eis;
    eis.element_index = ed.element_index;
    for (const auto& gp : ed.gauss_points) {
        eis.gauss_point_stresses.push_back(gp.stress);
    }
    bcs.add_initial_stress(std::move(eis));
}
\end{lstlisting}

\noindent
For this to take effect, the element assembly must incorporate
$\bsigma_0$ into the internal force:
\[
  \mathbf{f}_{\text{int}}(\mathbf{u})
  = \int_{\Omega_e} \mathbf{B}^T \bigl(\bsigma(\bvarepsilon(\mathbf{u}))
      + \bsigma_0\bigr)\,\dd V
\]
This requires modifications to \texttt{ContinuumElement::compute\_internal\-\_force\_vector},
specifically adding a \texttt{$\sigma$\_initial} term at each Gauss point.
This integration is deferred to a future phase.


% ═════════════════════════════════════════════════════════════════════════
\section{Comparison with OpenSees and Kratos}
\label{sec:bc-comparison}

\begin{table}[ht]
\centering\small
\caption{BC/IC handling comparison across frameworks.}
\label{tab:bc-framework}
\begin{tabular}{p{3cm}p{3.5cm}p{3.5cm}p{3.5cm}}
\toprule
\textbf{Feature} & \textbf{OpenSees} & \textbf{Kratos} & \textbf{fall\_n} \\
\midrule
Single-DOF fix
  & \texttt{fix nodeTag dof}
  & \texttt{DISPLACEMENT\_X: 0.0}
  & \texttt{constrain\_dof(n,d,v)} \\
Geometric selection
  & Tcl loops
  & Python \texttt{FindNodalH}
  & \texttt{PlaneSelector}, lambda \\
Non-zero Dirichlet
  & \texttt{imposedMotion} pattern
  & \texttt{ApplyDirichlet\-Process}
  & \texttt{constrain\_dof\_where(s,d,v)} \\
State transfer
  & \texttt{-database} serialise
  & \texttt{ModelPartIO}
  & \texttt{ModelState} capture/apply \\
Initial stress
  & Per-element update
  & \texttt{INITIAL\_STRESS\_TENSOR}
  & Defined, not yet integrated \\
Dispatch cost
  & Virtual (runtime)
  & Virtual + Python (runtime)
  & Compile-time (zero) \\
\bottomrule
\end{tabular}
\end{table}

\paragraph{Key differences:}
\begin{itemize}[nosep]
  \item \textbf{OpenSees} uses a tag-based Tcl scripting layer for BC
        specification.  Geometric loops require user scripting.  State transfer
        uses a persistence layer (\texttt{FE\_Datastore}).
  \item \textbf{Kratos} uses Python processes (\texttt{ApplyDirichletProcess},
        \texttt{ApplyNeumannProcess}) that iterate over model parts.
        Highly flexible but with Python $\to$ C++ boundary overhead.
  \item \textbf{fall\_n} resolves all selector logic at compile time.
        The constraint map lives in pure C++ with PETSc-section-based
        enforcement.  No scripting layer, no virtual dispatch per node.
\end{itemize}


% ═════════════════════════════════════════════════════════════════════════
\section{Extensibility Roadmap}
\label{sec:bc-extensibility}

The current infrastructure is designed for forward extension.
\Cref{tab:bc-roadmap} outlines planned capabilities and their
implementation path.

\begin{table}[ht]
\centering\small
\caption{Extensibility roadmap.}
\label{tab:bc-roadmap}
\begin{tabular}{p{3.5cm}p{5cm}p{5cm}}
\toprule
\textbf{Capability} & \textbf{Implementation path} & \textbf{Status} \\
\midrule
Per-DOF constraints
  & \texttt{constrain\_dof}, \texttt{constrain\_dof\_where}
  & \textbf{Implemented} \\
Non-zero Dirichlet
  & \texttt{fill\_imposed\_solution()} + NL correction
  & \textbf{Implemented} (NL only) \\
Node selectors
  & \texttt{NodeSelector} concept + built-in selectors
  & \textbf{Implemented} \\
State capture/transfer
  & \texttt{ModelState} + \texttt{capture\_state}
  & \textbf{Implemented} \\
\addlinespace
Initial stress
  & Extend \texttt{ContinuumElement::compute\_internal\_force}
    to add $\bsigma_0$ at each GP
  & Designed, deferred \\
Body forces
  & Consistent nodal load via $\int N^T\,\mathbf{b}\,\dd V$
  & Planned \\
Temperature field
  & Additional \texttt{ModelState} field + thermal strain
    $\bvarepsilon_\theta = \alpha\,\Delta T\,\mathbf{1}$
  & Planned \\
Multi-step builder
  & Fluent builder with chained \texttt{.constrain().material().build()}
  & Planned \\
State serialisation
  & HDF5/JSON export of \texttt{ModelState}
  & Planned \\
Linear non-zero BC
  & RHS correction $\mathbf{f}_{\text{eff}} = \mathbf{f}_{\text{ext}}
    - \mathbf{K}_{fc}\,\mathbf{u}_D$  in \texttt{LinearAnalysis}
  & Planned \\
\bottomrule
\end{tabular}
\end{table}


% ═════════════════════════════════════════════════════════════════════════
\section{Test Suite}
\label{sec:bc-tests}

The implementation is validated by 81 tests in
\texttt{tests/test\_boundary\_conditions.cpp}, organized as follows:

\begin{table}[ht]
\centering\small
\caption{Test suite structure (81 tests).}
\label{tab:bc-tests}
\begin{tabular}{lrl}
\toprule
\textbf{Category} & \textbf{Count} & \textbf{Coverage} \\
\midrule
Node selectors (pure geometry) & 38 &
  PlaneSelector, BoxSelector, SphereSelector, NodeIdSelector,
  combinators, lambdas, 2D/3D, factories \\
ModelState data structure & 9 &
  Default state, add/query, copy semantics \\
Per-DOF constraints & 11 &
  Single DOF, multiple DOFs, non-zero value, value update \\
Selector-based constraints & 10 &
  PlaneSelector on DOF, all DOFs, values, lambda, composite \\
Full solve: roller vs.\ clamped & 3 &
  Clamped $>$ 0, Roller $>$ 0, Roller $>$ Clamped \\
State capture after solve & 7 &
  Non-empty, node count, element count, non-zero
  displacement, non-zero stress, ID match \\
Backward compatibility & 3 &
  \texttt{fix\_x} gives positive deflection, reasonable magnitude \\
\bottomrule
\end{tabular}
\end{table}


% ═════════════════════════════════════════════════════════════════════════
\section{Summary}
\label{sec:bc-summary}

The generic BC/IC infrastructure extends \texttt{fall\_n} from an
all-or-nothing clamped-support tool to a framework capable of:

\begin{itemize}[nosep]
  \item \textbf{Arbitrary Dirichlet BCs:}  per-DOF, per-node, per-selector,
        with prescribed values (zero or non-zero).
  \item \textbf{Zero-cost selectors:}  compile-time geometric predicates
        with no runtime overhead versus hand-written loops.
  \item \textbf{State snapshots:}  PETSc-independent \texttt{ModelState}
        for continuation, restart, and future serialisation.
  \item \textbf{Full backward compatibility:}  \texttt{fix\_x}, \texttt{fix\_y},
        \texttt{fix\_z} continue to work exactly as before.
\end{itemize}

\noindent
The design follows the same policy-based, compile-time polymorphic
philosophy established in the element architecture (\Cref{ch:architecture})
and validated by an 81-test suite.


% ─── Bibliography ────────────────────────────────────────────────────────
\bibliographystyle{plain}
\bibliography{references}

\end{document}
